\documentclass[11pt]{article}

\usepackage{mathtools}
\usepackage{amssymb}
\usepackage{amsthm}
\usepackage{hyperref}
\usepackage{bm}
\usepackage{booktabs}
\usepackage[ruled,linesnumbered]{algorithm2e}
\usepackage{svg}
\usepackage{listings}
\usepackage{multirow}
\usepackage[table]{xcolor}
\usepackage{siunitx}
\usepackage{todonotes}

\usepackage{geometry}
\geometry{
	a4paper,
	margin=25mm
}

\lstdefinestyle{Python}{
  language=Python,
  basicstyle=\ttfamily\small,
  keywordstyle=\bfseries,
  commentstyle=\color{gray},
  numbers=left,
  numberstyle=\scriptsize\textbf,
  numbersep=4pt,
  xleftmargin=0pt,
  mathescape
}

\graphicspath{{thesis-media}}

\SetKwInput{KwReq}{Require}
\SetKwComment{Comment}{$\triangleright$\ }{}
\SetKwFor{ParFor}{parallel for}{do}{end parallel for}
\SetKwProg{Function}{function}{}{end}
\newcommand*{\Break}{\textbf{break}}
\newcommand{\Call}[2]{\operatorname{#1}\!\left(#2\right)}
\SetCommentSty{textrm}
\DontPrintSemicolon

\setcounter{MaxMatrixCols}{20}

\newcommand*{\closure}[1]{
	\ensuremath{\overline{\raisebox{0pt}[1.1\height]{$#1$}}}
}
\newcommand*{\slice}{\ensuremath{\, . \, . \,}}

\graphicspath{{images}}

\usepackage{xcolor}
\hypersetup{
	colorlinks,
	linkcolor={red!50!black},
	citecolor={blue!50!black},
	urlcolor={blue!80!black},
	linktocpage=true
}

\usepackage[sorting=nyt, style=ieee, backend=biber, doi=false, isbn=false,
  eprint=false, url=false, citestyle=numeric-comp]{biblatex}
\addbibresource{bibliography.bib}
\appto\bibfont{\setlength{\emergencystretch}{.5em}}

\title{A Nonoverlapping Additive Schwarz Method on a GPU}

% Tu jest dobre miejsce na Twoje własne makra i~środowiska:
\theoremstyle{definition}
\newtheorem{definition}{Definition}
\newtheorem{remark}{Remark}
\newtheorem*{remark*}{Remark}
\theoremstyle{theorem}
\newtheorem{theorem}{Theorem}
\newtheorem{conjecture}{Conjecture}
\newtheorem{lemma}{Lemma}
\newtheorem{corollary}{Corollary}
\newtheorem*{corollary*}{Corollary}

% koniec definicji

\begin{document}

\maketitle

% \tableofcontents
%\listoffigures
%\listoftables


\section{Introduction}

The objective of this thesis is to develop an efficient solver targeting
Graphics Processing Units (GPUs) \cite{cuda-guide} for the linear systems
arising from the discretization of elliptic problems of the form
\begin{equation}
  \begin{array}{rl}
    -\operatorname{div}(\rho \nabla u)=f\quad & \text{in }\Omega,\\
    u=0\quad & \text{on }\partial \Omega,
  \end{array}
  \label{div_eq}
\end{equation}
which occur, for example, in steady-state heat conduction~\cite{Ern-FEM}.
Here $f$ is a prescribed source term, $\rho$ is a given medium property
(e.g. conductivity) that may have jumps in heterogeneous media,
and $u$ is the unknown (e.g. temperature).

The number of unknowns in these systems can be extremely large, so exploiting the
coefficient matrix's special properties is necessary.
Specialized direct sparse solvers -- which perform a triangular
decomposition -- cannot handle the largest systems. Therefore, iterative methods
such as the Conjugate Gradient (CG), which compute only approximate solutions,
must be used instead \cite{Saad}.
Their convergence rate, in turn, depends on the system's condition number, which
typically grows fast with the number of unknowns~\cite{Ern-FEM};
effective preconditioners are therefore required to make an iterative method robust.
While general-purpose preconditioners such as Jacobi, Gauss-Seidel,
Incomplete LU or Incomplete Cholesky are available
\cite{Saad, cusparse}, problem-specific approaches often deliver superior
performance.

A broad class of such tailored preconditioners comprises domain decomposition
methods~\cite{toselli}, which partition the domain $\Omega$ into smaller subdomains,
solve local problems independently and then combine the results. In this thesis we
implement the three-grid, two-level Additive Schwarz Method (ASM) proposed and
analyzed in \cite{Przykry, Przykry-hp}. Designed as a preconditioner for the Discontinuous Galerkin
(DG) discretization of \eqref{div_eq} \cite{Ern-dg,Ern-weighted-dg}, it is robust
with respect to discretization parameters even in the presence of jumps in~$\rho$.
The DG discretization permits a~nonoverlapping ASM, which reduces coupling between
local problems. This is particularly important for distributed-memory systems,
since it minimizes inter-node communication, but is also beneficial
for shared-memory architectures such as GPUs, as it reduces data dependencies
and enables more parallelism.
Moreover, the preconditioner
construction allows a controlled balance between parallel scalability and
preconditioner quality, including massively parallel variants potentially
well-suited to GPUs~\cite{cuda-guide}.

Domain decomposition methods have been successfully implemented on
GPUs before. However, the amount of
work in this area seems marginal given the growing importance of GPUs in
scientific computing. Existing efforts typically handle
tens~\cite{Yamazaki,Sistek} or at most thousands~\cite{Papadrakakis} of subdomains,
primarily because each thread handles one subdomain at a~time.
They consider hybrid CPU-GPU implementations~\cite{Yamazaki, Sistek, Papadrakakis},
usually in a multi-node setting~\cite{Yamazaki,Sistek}, often porting
existing CPU software to GPUs~\cite{Yamazaki}.

Our approach differs in two key aspects. First, we target moderately sized problems
that fit on a single GPU. We design our solver from the ground up for a GPU
so that CPU does not participate in the solve phase and is used only for
orchestration. Second, we implement the three-grid ASM~\cite{Przykry, Przykry-hp},
which, to our knowledge, has not been ported to a GPU before. This allows
us to consider highly parallel variants with millions of small subdomains,
processing all of them simultaneously within a single GPU kernel; this design seems
better suited to the GPU architecture \cite{cuda-guide}.\todo{coś o gpu vs cpu?}
\todo{coś o pytorchu?}

We assume that the linear system $\bm{A} \bm{u^*} = \bm{f}$, resulting from the
DG discretization of~\eqref{div_eq}, has already been assembled and is provided.
Our first goal is to implement an efficient single\nobreakdash-GPU solver for this system,
employing the preconditioned CG method accelerated by the ASM preconditioner.
Our second goal is to evaluate its performance across a range of parameters
to identify the best-performing configurations and to compare it with the
standard two-grid variant~\cite{classical-estimate}.
Finally, we will compare the preconditioner and its implementation
against NVIDIA's state-of-the-art GPU preconditioner -- AmgX \cite{amgx}.

The thesis is organized as follows. Chapter~\ref{chap:theory} presents
the mathematical foundations, defines the ASM preconditioner
and derives its algorithmic realization. Chapter~\ref{chap:implementation}
details the implementation of the preconditioner and other solver components,
including the CG method. Finally, Chapter~\ref{chap:experiments} reports
experimental results across a range of parameter settings, identifies
the best-performing configurations and compares them with the
NVIDIA AmgX library.

\section{Numerical Framework}
\label{chap:theory}

\todo{nieaktualne} The objective of this chapter is to present the mathematical foundations required for
understanding and, henceforth, implementing an Additive Schwarz Method (ASM) solver for \eqref{div_eq}.
Section~\ref{sec:variational-formulation} introduces the variational formulation
of \eqref{div_eq}, which serves as the basis for the
Discontinuous Galerkin (DG) discretization described in Section~\ref{sec:discretization}.
In Section~\ref{sec:additive-schwarz-method}, we define
the ASM preconditioner for the obtained discrete system and examine its key properties. Finally,
Section~\ref{sec:preconditioner-alg} derives the detailed algorithmic procedure
for applying the ASM preconditioner, which provides the foundation for the
GPU\nobreakdash-accelerated implementation presented in Chapter~\ref{chap:implementation}.

\subsection{Variational Formulation of the Problem}
\label{sec:variational-formulation}

Let $\Omega \subset \mathbb{R}^d$, where $d\in \{2, 3\}$, be an open, bounded
polyhedral domain with Lipschitz boundary $\partial \Omega$. We denote by
$(\cdot, \cdot)_{L^2(\Omega)}$ the inner product on $L^2(\Omega)$ and by
$\| \cdot \|_{L^2(\Omega)}$ the corresponding norm:
\[
  (u, v)_{L^2(\Omega)} = \int_\Omega u(x) v(x)\, dx,\quad
  \| u \|_{L^2(\Omega)} = (u, u)_{L^2(\Omega)}^{1/2},\quad
  \forall u, v\in L^2(\Omega).
\]
Similarly, we denote the norm and semi-norm on the Sobolev space $H^k(\Omega)$
by $\|\cdot\|_{H^k(\Omega)}$ and $|\cdot|_{H^k(\Omega)}$, respectively,
defined by
\begin{equation*}
  \|u\|_{H^k(\Omega)} = \left(\sum_{|\alpha| \leq k} \|D^\alpha u\|_{L^2(\Omega)}^2\right)^{1/2},\quad
  |u|_{H^k(\Omega)} = \left(\sum_{|\alpha| = k} \|D^\alpha u\|_{L^2(\Omega)}^2 \right)^{1/2}.
\end{equation*}

We assume that
${f\in L^2(\Omega)}$, ${\rho\in L^\infty(\Omega)}$ are prescribed, and that there
exist positive constants $c_0, c_1$ such that $c_0 \leq \rho \leq c_1$ almost everywhere
in $\Omega$. Let the bilinear form
${a: H^1(\Omega)\times H^1(\Omega) \to \mathbb{R}}$ be given by
\[
  a(u, v) = \int_\Omega \rho(x) \nabla u(x) \cdot \nabla v(x)\, dx.
\]
The \emph{variational formulation} of \eqref{div_eq} is then:
Find $u^* \in H_0^1(\Omega)$ such that
\begin{equation}
  a(u^*, v) = (f, v)_{L^2(\Omega)} \quad \forall v \in H_0^1(\Omega).\label{var-cont-formulation}
\end{equation}
Owing to the boundedness of $\rho$, the bilinear form $a$ is continuous and coercive
on $H_0^1(\Omega)$, which guarantees existence and uniqueness of the solution
by the Lax-Milgram lemma \cite{Brenner}.

\subsection{Discontinuous Galerkin Discretization of the Problem}
\label{sec:discretization}

Let $\{\mathcal{T}_h\}_{h>0}$ be a shape-regular family of affine
triangulations of $\Omega$ (cf. \cite[Definitions 1.53, 1.107]{Ern-FEM}), and denote by
$h=\max_{K\in \mathcal{T}_h} \operatorname{diam}(K)$ a parameter describing the mesh size.
We fix one triangulation $\mathcal{T}_h$ (the \emph{fine mesh}) and
denote by $\mathcal{E}_h$ the set of all its faces (edges in 2D).
For a~given integer $p\geq 1$, we will approximate the solution of \eqref{var-cont-formulation}
in the \emph{broken finite element space}
\[
  V_h = \{v \in L^2(\Omega) : v|_K \in \mathcal{P}_p \text{ for all } K\in \mathcal{T}_h\},
\]
where $\mathcal{P}_p$ is the space of real polynomials of degree at most $p$.
For $u, v\in V_h$, $K\in \mathcal{T}_h$, $e\in \mathcal{E}_h$ we use
the standard notation
\[
  (u, v)_K = \int_K u(x) v(x)\, dx,\quad \langle u, v \rangle_e = \int_e u(x) v(x)\, d\sigma_e(x),
\]
where $\sigma_e$ is the surface measure on $e$. Discretization of
\eqref{var-cont-formulation} by the symmetric weighted interior penalty
Discontinuous Galerkin (DG) method leads to the following problem \cite{Przykry, Ern-weighted-dg}:
Find $u_h^* \in V_h$ such that
\begin{equation}
  \mathcal{A}_h(u_h^*, v_h) = (f, v_h)_{L^2(\Omega)} \quad \forall v_h \in V_h,\label{var-formulation}
\end{equation}
where $\mathcal{A}_h: V_h \times V_h \to \mathbb{R}$ is a bilinear form on $V_h$ defined as
\begin{equation}
\mathcal{A}_h(u, v) = \sum_{K\in \mathcal{T}_h} \left(\rho \nabla u, \nabla v\right)_K
  + \sum_{e\in \mathcal{E}_h} \langle \gamma [u], [v]\rangle_e
  - \sum_{e\in \mathcal{E}_h} \langle [u], \{\rho \nabla v\}_\omega\rangle_e
  - \sum_{e\in \mathcal{E}_h} \langle \{\rho \nabla u\}_\omega, [v] \rangle_e.
  \label{A_h_def}
\end{equation}
For an interior face $e=\partial K_1 \cap \partial K_2$, $K_1, K_2 \in \mathcal{T}_h$, we set
\[
  \gamma = \frac{\delta p^2}{h_e}\frac{\rho_1 \rho_2}{\rho_1 + \rho_2},\quad
  \omega_1 = \frac{\rho_2}{\rho_1 + \rho_2},\quad
  \omega_2 = \frac{\rho_1}{\rho_1 + \rho_2},
\]
where $\rho_i=\rho|_{K_i}$, $h_e=\operatorname{diam}(e)$ and $\delta>0$ is a prescribed constant.
The \emph{weighted average} and \emph{jump} operators are then
\[
  \{\rho \nabla u\}_{\omega} = \omega_1 \rho_1 \nabla u_1 + \omega_2 \rho_2 \nabla u_2,\quad
  [u] = u_1 n_1 + u_2 n_2,
\]
with $u_i=u|_{K_i}$, and $n_i$ being the outward normal
vector to $K_i$ on $e$ for $i=1,2$. On a boundary face $e\subseteq \partial K_1\cap \partial \Omega$,
we take
\[
  \gamma = \frac{\delta p^2}{h_e}\rho_1,\quad
  \{\rho \nabla u\}_\omega = \rho_1 \nabla u_1,\quad
  [u] = u_1 n_1.
\]
Values of $\{\rho \nabla v\}$ and $[v]$ are defined analogously.

Let us fix a basis $\Phi=(\varphi_0, ..., \varphi_{n-1})$ of the finite element space $V_h$.
Elements of the \emph{stiffness matrix} ${\bm{A}=\left[a_{ij}\right]}$ are defined
as $a_{ij} = \mathcal{A}_h(\varphi_i, \varphi_j)$ for $i,j=0,...,n-1$. Equivalently,
$\bm{A}$ is the matrix representation of the bilinear form $\mathcal{A}_h$
with respect to the basis $\Phi$.
Similarly, let us define elements of the \emph{load vector} $\bm{f}=(f_0, ..., f_{n - 1})$
as $f_i = (f, \varphi_i)_{L^2(\Omega)}$. The variational problem
\eqref{var-formulation} is then equivalent to the linear system
\begin{equation}
  \bm{A} \bm{u^*} = \bm{f},\label{auf}
\end{equation}
where $\bm{u^*}=(u_0, ..., u_{n - 1})\in\mathbb{R}^n$ are the coefficients
of $u_h^*$ in the basis $\Phi$, i.e. $u_h^*=\sum_{i=0}^{n-1} u_i \varphi_i$.
In order for both \eqref{var-formulation} and \eqref{auf} to
be well-posed, the matrix $\bm{A}$ must
be invertible, which is guaranteed by the following result.

\begin{theorem}[\cite{Ern-weighted-dg}]\label{thm:spd}
  There exists $\delta_0>0$ such that for all $h>0$, $p\geq 1$ and $\delta \geq \delta_0$,
  the bilinear form~$\mathcal{A}_h$ (and hence its corresponding matrix $\bm{A}$)
  is symmetric and positive definite.
\end{theorem}

From this point onward, we shall assume that $\delta \geq \delta_0$, and consequently,
the conclusions of the above theorem hold. Under this assumption, the
following approximation result applies.

\begin{theorem}[\cite{Ern-weighted-dg}]
  If the exact solution $u^*\in H_0^1(\Omega)$ of
  \eqref{var-cont-formulation} belongs to the broken Sobolev space
  \[
    H^{p+1}(\mathcal{T}_h)=\{v \in L^2(\Omega) : v|_K \in H^{p+1}(K) \text{ for all } K\in \mathcal{T}_h\},
  \]
  then there exists a constant $C>0$ independent of $h$ such that
  \[
    \|u_h^* - u^*\|_{L^2(\Omega)} \leq C h^{p+1} \|u^*\|_{H^{p+1}(\mathcal{T}_h)}.
  \]
\end{theorem}

\subsection{Schwarz Preconditioners for the Discrete System}
\label{sec:additive-schwarz-method}

We aim to solve the linear system \eqref{auf} using the Preconditioned
Conjugate Gradient (PCG) method \cite[Section 9.2]{Saad}, which requires an
effective preconditioner for the matrix $\bm{A}$. To this end, we employ
the three-grid nonoverlapping Schwarz Method from~\cite{Przykry},
which is tailored specifically for the system~\eqref{auf}.
We now outline this construction.

For this purpose we introduce two additional
partitions of~$\Omega$: \emph{local solvers' mesh} $\mathcal{T}_H=\{\Omega_0, ..., \Omega_{N_H-1}\}$ and
\emph{coarse mesh} $\mathcal{T}_{\mathcal{H}}=\{D_0, ..., D_{N_{\mathcal{H}-1}}\}$.
We assume the inclusion
\[
  \mathcal{T}_{\mathcal{H}}\subseteq \mathcal{T}_H\subseteq \mathcal{T}_h,
\]
in the sense that $\mathcal{T}_H$ is a refinement of $\mathcal{T}_{\mathcal{H}}$ and
$\mathcal{T}_h$ is a refinement of $\mathcal{T}_H$ (see Figure~\ref{fig:mesh_examples}
for examples). Finally, set
$H=\max_{j} \operatorname{diam}(\Omega_j)$ and $\mathcal{H}=\max_{j} \operatorname{diam}(D_j)$.

Following \cite{Przykry}, we decompose the finite element space $V_h$, hereafter
denoted simply by $V$, into the sum
\begin{equation}
  V = V_C + \sum_{i=0}^{N_H-1} V_i,\label{V_decomp}
\end{equation}
where the \emph{coarse space} is
\[
  V_C = \{v \in V : v|_{D_j}\in \mathcal{P}_0\text{ for } j=0,...,N_{\mathcal{H}}-1\}
\]
and the \emph{local spaces} are
\[
  V_i = \{v\in V : v|_{\Omega_j}=0 \text{ for all } j\neq i\},\quad i=0,...,N_H-1.
\]
Note that $V$ is a direct sum of the local spaces $V_0, ..., V_{N_H-1}$.

We define the \emph{coarse solve operator}
$T_C: V\to V_C$ by the relation
\begin{equation}
  \mathcal{A}_h(T_C u, v_C) = \mathcal{A}_h(u, v_C) \quad \forall v_C\in V_C,\label{T_C_def}
\end{equation}
and analogously the \emph{local solve operators} $T_i: V\to V_i$ by
\begin{equation}
  \mathcal{A}_h(T_i u, v_i) = \mathcal{A}_h(u, v_i) \quad \forall v_i\in V_i\label{T_i_def}
\end{equation}
for $i=0,...,N_H-1$. The preconditioned additive Schwarz operator $T_{add}:V\to V$
corresponding to \eqref{V_decomp} is then defined as
\begin{equation}
  T_{add} = T_C + \sum_{i=0}^{N_H-1} T_i.\label{T_def}
\end{equation}
In addition, we consider
the preconditioned hybrid Schwarz operator $T_{hyb}:V\to V$, defined as
\begin{equation}
  T_{hyb} = T_C + \left(I - T_C\right) \sum_{i=0}^{N_H-1} T_i(I-T_C).\label{T_hyb_def}
\end{equation}

Note that, unlike in \cite{Przykry}, we do not employ
a simplified, inexact form in the definitions of $T_C$ and $T_i$.
The simplified form can be beneficial in a matrix\nobreakdash-free
context, where its values are evaluated on-the-fly. However, it yields no
computational advantage if assembled in advance,
as its matrix has the same sparsity
pattern as the original form $\mathcal{A}_h$. Furthermore, maintaining both forms
would double the memory footprint, significantly reducing size of the problem
that can be solved on a single GPU.

To utilize the PCG method, the operators $T_{add}$ and $T_{hyb}$ must be
symmetric and positive definite (SPD) with respect to the inner product induced
by $\mathcal{A}_h$. Sufficient conditions for $T_{add}$ to satisfy these properties,
along with estimates of its condition number, are established in \cite{Przykry}
and \cite{Przykry-hp}. These results can be extended to $T_{hyb}$ using tools
from \cite{toselli}. Since the implementation is independent of the specific
details, we omit them and proceed under the assumption that $\rho$, $\mathcal{T}_h$,
$\mathcal{T}_H$ and $\mathcal{T}_{\mathcal{H}}$ are chosen such that
both $T_{add}$ and $T_{hyb}$ are SPD.

We now derive a formula for the preconditioning matrices $\bm{M_{add}}$ and
$\bm{M_{hyb}}$, which follow from the definition of the preconditioned operators
$T_{add}$ and $T_{hyb}$, respectively.
Assume that $\Phi_C$ is a basis of $V_C$ and that $\Phi_i$ is a~basis of $V_i$ for
$i=0,...,N_H-1$. Whenever we introduce matrices of
linear operators defined on the spaces $V$, $V_C$ and $V_i$ for $i=0,...,N_H-1$, we will
always assume they are represented with respect to $\Phi$, $\Phi_C$ and $\Phi_i$,
respectively, for $i=0,...,N_H-1$. The specific bases $\Phi$, $\Phi_C$ and $\Phi_i$ will later
be chosen to facilitate an efficient implementation of the preconditioner.

Let $\bm{T_C}$, $\bm{T_i}$, $\bm{T_{add}}$ and $\bm{T_{hyb}}$ be the matrices of the
operators $T_C$, $T_i$, $T_{add}$ and $T_{hyb}$, respectively, for ${i=0,...,{N_H-1}}$.
Let $\bm{R_C}^T$ denote the matrix of the identity embedding $R_C^T: V_C \to V$
and let $\bm{R_C}$ be its transpose. Matrices
$\bm{R_i}^T$ and $\bm{R_i}$ are defined analogously for the local spaces $V_i$,
$i=0,...,N_H-1$. From \eqref{T_C_def} we obtain
\[
\left(\bm{R_C}^T\bm{v_C}\right)^T \bm{A} \left(\bm{R_C}^T\bm{T_C}\bm{u}\right)
  = \left(\bm{R_C}^T\bm{v_C}\right)^T \bm{A} \bm{u}
\quad \forall \bm{u}\in \mathbb{R}^n\ \forall \bm{v_C}\in \mathbb{R}^{\dim V_C},
\]
which simplifies to
\begin{equation}
\bm{R_C} \bm{A} \bm{R_C}^T \bm{T_C} = \bm{R_C} \bm{A},\quad i=0,...,N_H.\label{T_C_eq}
\end{equation}
Since $\bm{A}$ is SPD, and matrices $\bm{R_C}$ and $\bm{R_C}^T$ have full rank
(as $R_C^T:V_C\to V$ is an embedding), the matrix
\begin{equation}
\bm{A_C} = \bm{R_C} \bm{A} \bm{R_C}^T\label{A_C_def}
\end{equation}
is also SPD. In particular, it is invertible, and
\eqref{T_C_eq} can be rewritten as
\begin{equation}
\bm{T_C} = \bm{A_C}^{-1}\bm{R_C}\bm{A}.\label{T_C_matrix}
\end{equation}
Analogously, from \eqref{T_i_def} we obtain
\begin{equation}
\bm{T_i} = \bm{A_i}^{-1}\bm{R_i}\bm{A},\quad i=0,...,N_H-1,\label{T_i_matrix}
\end{equation}
where
\begin{equation}
\bm{A_i} = \bm{R_i} \bm{A} \bm{R_i}^T,\quad i = 0, ..., N_H - 1.\label{A_i_def}
\end{equation}
The matrix $\bm{A_C}$ will be referred to as the \emph{coarse solver matrix},
while $\bm{A_0}, ..., \bm{A_{N_H-1}}$ will be called the \emph{local solvers' matrices}.

Now, from \eqref{T_def}, we obtain
\[
\bm{T_{add}} = \bm{R_C}^T \bm{T_C} + \sum_{i=0}^{N_H-1}\bm{R_i}^T\bm{T_i} =
  \bm{R_C}^T\bm{A_C}^{-1}\bm{R_C}\bm{A}
  + \sum_{i=0}^{N_H-1} \bm{R_i}^T\bm{A_i}^{-1}\bm{R_i}\bm{A}=:\bm{M_{add}}\bm{A},
\]
where
\begin{equation}
  \bm{M_{add}} = \bm{R_C}^T \bm{A_C}^{-1} \bm{R_C} + \sum_{i=0}^{N_H-1} \bm{R_i}^T \bm{A_i}^{-1} \bm{R_i},\label{M_matrix}
\end{equation}
is the additive Schwarz preconditioner for the system \eqref{auf}. Similarly,
by combining \eqref{T_hyb_def}, \eqref{T_C_matrix} and \eqref{T_i_matrix}, we obtain
\begin{multline*}
  \bm{T_{hyb}}= \bm{R_C}^T \bm{T_C} + \left(\bm{I} - \bm{R_C}^T \bm{T_C}\right)
  \sum_{i=0}^{N_H-1} \bm{R_i}^T \bm{T_i} \left(\bm{I} - \bm{R_C}^T \bm{T_C}\right)=\\
  = \bm{R_C}^T \bm{A_C}^{-1} \bm{R_C} \bm{A} +
  \left(\bm{I} - \bm{R_C}^T \bm{A_C}^{-1} \bm{R_C} \bm{A}\right)
  \sum_{i=0}^{N_H-1} \bm{R_i}^T \bm{A_i}^{-1} \bm{R_i}
  \left(\bm{I} - \bm{A}\bm{R_C}^T \bm{A_C}^{-1} \bm{R_C}\right)\bm{A}=: \bm{M_{hyb}}\bm{A}.\\
\end{multline*}
Here we set
\[
  \bm{M_{hyb}} = \bm{R_C}^T \bm{A_C}^{-1} \bm{R_C} + \bm{M_{hyb}^{loc}},
\]
where
\begin{equation}
  \bm{M_{hyb}^{loc}} =
  \left(\bm{I} - \bm{R_C}^T \bm{A_C}^{-1} \bm{R_C} \bm{A}\right)
  \sum_{i=0}^{N_H-1} \bm{R_i}^T \bm{A_i}^{-1} \bm{R_i}
  \left(\bm{I} - \bm{A}\bm{R_C}^T \bm{A_C}^{-1} \bm{R_C}\right).
  \label{M_hyb_loc}
\end{equation}
In contrast to the additive case, we do not solve the preconditioned system
\[
  \bm{M_{hyb}} \bm{A} \bm{u^*} = \bm{M_{hyb}} \bm{f}
\]
directly using PCG. Instead, we employ the \emph{projected conjugate gradient} algorithm
\cite{toselli}. To this end, we set
\[
  \bm{u_C} = \bm{R_C^T}\bm{A_C}^{-1}\bm{R_C}\bm{A} \bm{u^*}= \bm{R_C^T} \bm{A_C}^{-1} \bm{R_C} \bm{f}
\]
and observe that
\[
  \bm{M_{hyb}} \bm{A} \bm{u^*} = \bm{u_C} + \bm{M_{hyb}^{loc}} \bm{A} \bm{u^*}=
  \bm{u_C} + \bm{M_{hyb}^{loc}} \bm{A} \left(\bm{u^*} - \bm{u_C}\right).
\]
Therefore, we may equivalently solve
\begin{equation}
  \bm{M_{hyb}^{loc}} \bm{A} \bm{v} =
  \bm{M_{hyb}^{loc}} \bm{f}, \label{projected_system}
\end{equation}
for $\bm{v}\in \operatorname{range}(\bm{I} - \bm{R_C^T} \bm{A_C}^{-1} \bm{R_C} \bm{A})$
and recover the solution of the original system as $\bm{u^*} = \bm{u_C} + \bm{v}$.

\subsection{Algorithmic Realization of the Schwarz Preconditioners}
\label{sec:preconditioner-alg}

The PCG method uses the preconditioning matrices $\bm{M_{add}}$ and $\bm{M_{hyb}^{loc}}$
defined in \eqref{M_matrix} and \eqref{M_hyb_loc}, respectively,
only to compute their product with the residual vector at each iteration.
Therefore we do not need to assemble $\bm{M_{add}}$ or $\bm{M_{hyb}^{loc}}$
explicitly, but just to efficiently evaluate the matrix\nobreakdash-vector product.
In this section we present parallel algorithms for that purpose.

Here and in following implementation, we assume that the number of degrees of freedom
(DoFs) in each local space $V_i$ is the same; that is $\dim V_i = k$ for
all $i=0,...,N_H-1$ and some fixed integer $k$.
For example, this condition holds
when $\mathcal{T}_h$ is a uniform refinement of $\mathcal{T}_H$,
which includes the massively parallel variant where each local problem is posed on
a~single element. Under this assumption, all local solvers' matrices
$\bm{A_i}$ have identical dimensions $k\times k$, which is crucial for the
efficient GPU parallelization (see Section~\ref{chap:implementation}). We do not,
however, require the matrices $\bm{A_i}$ to be identical, preserving flexibility
in the mesh design and allowing for variable coefficient $\rho$.

Let $l_i$ be the number of local solvers' subdomains $\Omega_j$ that
are contained in the element $D_i$ of the coarse mesh $\mathcal{T}_{\mathcal{H}}$, i.e.
\begin{equation}
  l_i = \left|\{j : \Omega_j \subseteq D_i\}\right|,\quad i=0,...,N_{\mathcal{H}-1}.\label{l_i_def}
\end{equation}
Note that we allow different values of $l_i$ for different $i$. To shorten the notation,
let us set
\begin{equation}
  \bar l_0 = 0;\quad \bar l_i = l_1 + ... + l_{i-1},\quad i=1,...,N_{\mathcal{H}}.\label{bar_l_i_def}
\end{equation}
Observe that $\bar l_{N_{\mathcal{H}}} = N_H$.

Now, as the basis $\Phi=(\varphi_0, ..., \varphi_{n-1})$ of $V$ let us choose the
standard finite element nodal basis (see, e.g., \cite{Brenner}). This choice
ensures that the stiffness matrix $\bm{A}$ is sparse.
After an appropriate permutation of the indices, we may additionally assume that
\begin{align*}
\operatorname{supp} \varphi_j \subseteq \closure{\Omega_i} &\iff ki \leq j <  k(i + 1),\\
\Omega_j \subseteq D_i &\iff \bar l_{i} \leq j < \bar l_{i + 1},
\end{align*}
which implies
\begin{equation}
\operatorname{supp} \varphi_j \subseteq \closure{D_i}\iff k\bar l_{i} \leq j < k \bar l_{i + 1}.\label{l_i_res}
\end{equation}
Then we may choose
\[
\Phi_i = \left(\varphi_{j} : ki \leq j <  k(i + 1)\right)
\]
as the basis of $V_i$ for $i=0,...,N_H-1$ and
\[
\Phi_C = \left(\varphi_{k\bar l_{i}} + ... + \varphi_{k \bar l_{i + 1}-1} :
  i=0,...,N_{\mathcal{H}}-1\right)
\]
as the basis of $V_C$. As a result, for $i=0,...,N_H-1$, the matrix
$\bm{R_i}$ has the following, very simple block structure
\newcommand\undermat[2]{%
  \makebox[0pt][l]{%
    $\smash{\kern-0.85\arraycolsep
      \underbrace{\phantom{\kern\arraycolsep\begin{matrix}#2\end{matrix}\kern0.85\arraycolsep}}_{\text{$#1$}}}$}#2}
\begin{equation}
\arraycolsep=8pt\def\arraystretch{1.5}
  \bm{R_i}=\left[\begin{array}{c|c|c|c|c|c|c|c|c}
      \undermat{i\text{ blocks}}{\bm{0_k} & \bm{0_k} & \dots & \bm{0_k}} &
      \bm{I_k} &
      \undermat{(N_H - i - 1)\text{ blocks}}{\bm{0_k} & \bm{0_k} & \dots & \bm{0_k}}\\
  \end{array}\right]\vspace{20pt}
  \label{Ri_structure}
\end{equation}
where $\bm{I_k}$ denotes the $k\times k$ identity matrix and $\bm{0_k}$
the zero matrix of the same size. Therefore
\begin{equation}
  \bm{R_i}^T(u_{ki}, ..., u_{k(i + 1) - 1})=
    (\underbrace{0, 0, ..., 0}_{ki}, u_{ki}, ..., u_{k(i + 1) - 1},
    \underbrace{0, 0, ..., 0}_{k(N_H - i - 1)}),\quad i=0,...,N_H-1\label{R_i_T_effect}
\end{equation}
and
\begin{equation}
  \bm{R_i}(u_0, ..., u_{n-1}) = (u_{ki}, ..., u_{k(i + 1) - 1}),\quad i=0,...,N_H-1.\label{R_i_effect}
\end{equation}
For the coarse space we have
\begin{equation}
  \bm{R_C}=\begin{bmatrix}
    \multicolumn{4}{c}{\smash[b]{\underbrace{\begin{matrix}1 & 1 & \cdots & 1\end{matrix}}_{kl_0}}} &   &   &  &   &  &   &   & &  \\
      &   &  &   & \multicolumn{4}{c}{\smash[b]{\underbrace{\begin{matrix}1 & 1 & \cdots & 1\end{matrix}}_{kl_1}}} &  &   &   &  &  \\
      &  &  &  &  &  &  &  & \ddots &  &  &  & \\
      &  &  &  &  &  &  &  &  & \multicolumn{4}{c}{\smash[b]{\underbrace{\begin{matrix}1 & 1 & \cdots & 1\end{matrix}}_{kl_{N_\mathcal{H}-1}}}}
  \end{bmatrix}\vspace{1 pt}\label{R_C_structure},
\end{equation}
hence
\begin{equation}
  \bm{R_C}^T(v_0, v_1, ..., v_{N_{\mathcal{H}}-1}) =
    (\underbrace{v_0, v_0, ..., v_0}_{kl_0},
    \underbrace{v_1, v_1, ..., v_1}_{kl_1}, ...,
    \underbrace{v_{N_{\mathcal{H}}-1}, v_{N_{\mathcal{H}}-1}, ..., v_{N_{\mathcal{H}}-1}}_{kl_{N_\mathcal{H}-1}}),\label{R_C_T_effect}
\end{equation}
and
\begin{equation}
  \bm{R_C}(u_0, ..., u_{n-1}) = (c_0, ..., c_{N_{\mathcal{H}}-1}),\quad \text{where } c_i = u_{k\bar l_i} + ... + u_{k\bar l_{i + 1} - 1}.
  \label{R_C_effect}
\end{equation}

From \eqref{A_i_def} and \eqref{Ri_structure} we see that if
$\bm{A}=\left[a_{rs}\right]_{r,s=0,...,{n-1}}$, then
\begin{equation}
  \bm{A_i} = \left[a_{rs}\right]_{r,s=ki, ..., k(i + 1) - 1},\quad i=0,...,N_H-1,\label{A_i_structure}
\end{equation}
i.e. $\bm{A_i}$ is a submatrix of $\bm{A}$ containing only the rows and columns
indexed by $ki, ..., k(i + 1)-1$. On the other hand,
from \eqref{A_C_def} and \eqref{R_C_structure} we obtain
\begin{equation}
  \bm{A_C}=\left[c_{ij}\right]_{i,j=0,...,N_{\mathcal{H}}-1},
  \text{ where } c_{ij} = \sum_{r=k\bar l_{i}}^{k \bar l_{i+1}-1}
  \ \sum_{s=k\bar l_{j}}^{k\bar l_{j+1}-1} a_{rs}.\label{A_C_structure}
\end{equation}

For the sake of completeness, we note that there is an alternative way to
construct the matrix~$\bm{A_C}$. From \eqref{A_C_def} we know that $\bm{A_C}$
represents the restriction of $\mathcal{A}_h$ to the coarse
space~$V_C$. Moreover, for $u_C, v_C\in V_C$ one has
$\nabla u_C \equiv \nabla v_C\equiv0$ almost everywhere in each $D_j$ for
$j=0,...,N_{\mathcal{H}}-1$, and hence \eqref{A_h_def} reduces to
\[
\mathcal{A}_h(u_C, v_C)= \sum_{e\in \mathcal{E}_h} \langle \gamma [u_C], [v_C]\rangle_e,
\quad \forall u_C, v_C\in V_C.
\]
Therefore, $\bm{A_C}$ can be efficiently assembled directly from the
right-hand side of the formula above, provided the
basis of $V_C$ is available to the discretization software. This approach,
however, was not considered in this work, as it couples the discretization phase
with the selection of the solver. In any case,
the assembly method does not affect the solution time, which was our primary
optimization target.

One final algorithmic detail remains. In the implementation of
$\bm{M_{hyb}^{loc}}$, we need to compute the action of $\bm{A}$ immediately
followed by the restriction $\bm{R_C}$. We can optimize this operation by
precomputing the matrix $\bm{A_R} = \bm{R_C}\bm{A}$. From \eqref{R_C_structure}
it follows that
\begin{equation}
  \bm{A_R} = \left[d_{ij}\right]_{i=0,...,N_{\mathcal{H}}-1; j=0,...,n-1},
  \text{ where } d_{ij} = \sum_{s=k\bar l_{i}}^{k\bar l_{i+1}-1} a_{sj}.
  \label{A_R_structure}
\end{equation}

We now present the algorithmic realization of the Schwarz preconditioners.
The setup phase, common to both preconditioners, is outlined in
Algorithm~\ref{preconditioner-setup-pseudocode}. The functions \textit{M\_{add}}
and \textit{M\_{hyb}\_{loc}}, which multiply the respective preconditioning matrices
by a vector, are defined in Algorithms~\ref{preconditioner-M-pseudocode}
and \ref{preconditioner-M-hyb-pseudocode}, respectively.
Since these algorithms operate solely on matrix representations of the operators,
the boldface notation is omitted for notational clarity.

Notably, \textit{M\_{hyb}\_{loc}} assumes that the
input vector belongs to the range of $\bm{I} - \bm{A}\bm{R_C^T} \bm{A_C}^{-1} \bm{R_C}$.
As shown in \cite{toselli}, this condition is naturally satisfied
when the function is used within the PCG method applied to the projected system
\eqref{projected_system}.
Since $\bm{I} - \bm{A}\bm{R_C^T} \bm{A_C}^{-1} \bm{R_C}$ is a~projection (which
can be verified by direct calculation), it need not be explicitly applied at
the beginning of \textit{M\_{hyb}\_{loc}}, saving computational time.

\begin{algorithm}
  % When editing this algorithm, go through all references and check whether
  % the line numbers are correct.
  \caption{Schwarz Preconditioners -- Setup}
  \label{preconditioner-setup-pseudocode}
  \KwIn{$A$, $k$, $l=(l_0, ..., l_{N_{\mathcal{H}}-1})$}
  \KwOut{$\bar l$, \textsc{CoarseSolve}, \textsc{LocalSolve}$[0\slice N_H-1]$}

  $\bar l \gets \Call{CumulativeSums}{l}$  \Comment*[r]{\eqref{bar_l_i_def}}
  \ParFor(\Comment*[f]{$A_i\gets R_iAR_i^T$}){$i=0, ..., N_H-1$} {
    $A_i\gets A[ki \slice k(i + 1)-1][ki \slice k(i + 1)-1]$  \Comment*[r]{\eqref{A_i_structure}}
    $\textsc{LocalSolve}[i] \gets \Call{InitLocalSolver}{A_i}$  \Comment*[r]{e.g. Cholesky factorization of $A_i$}
  }
  \ParFor(\Comment*[f]{$A_C\gets R_CAR_C^T$}){$i,j=0, ..., N_{\mathcal{H}}-1$} {
    $A_C[i][j]\gets \Call{Sum}{A[k\bar l_{i} \slice k\bar l_{i + 1} - 1][k\bar l_{j} \slice k\bar l_{j + 1}-1]}$  \Comment*[r]{\eqref{A_C_structure}}
  }
  $\textsc{CoarseSolve} \gets \Call{InitCoarseSolver}{A_C}$  \Comment*[r]{e.g. sparse LU factorization of $A_C$}
  \ParFor(\Comment*[f]{$A_R\gets R_CA$ (only for $M_{hyb}^{loc}$)}){$i=0, ..., N_{\mathcal{H}}-1;\ j=0,...,n - 1$} {
    $A_R[i][j] \gets \Call{Sum}{A[k\bar l_{i} \slice k\bar l_{i + 1} - 1][j]}$  \Comment*[r]{\eqref{A_R_structure}}
  }
\end{algorithm}

\begin{algorithm}
  % When editing this algorithm, go through all references and check whether
  % the line numbers are correct.
  \caption{Additive Schwarz Preconditioner -- Multiplication By a Vector}
  \label{preconditioner-M-pseudocode}
  \KwReq{$k$, $l$, $\bar l$, \textsc{CoarseSolve}, \textsc{LocalSolve}$[0\slice N_H-1]$}

  \Function(\Comment*[f]{computes $M_{add}x$}){M\_add($x$)} {
    \ParFor{$i=0, ..., N_H-1$} {
      $x^{(i)}\gets x[ki \slice k(i + 1)-1]$  \Comment*[r]{$x^{(i)}\gets R_ix$ \eqref{R_i_effect}}
      $y^{(i)}\gets \textsc{LocalSolve}[i]\!\left(x^{(i)}\right)$  \Comment*[r]{$y^{(i)}\gets A_i^{-1}x^{(i)}$}
      $y[ki \slice k(i + 1) - 1] \gets y^{(i)}$  \Comment*[r]{$y \gets R_0^Ty^{(0)}+ ...+R_{N_H-1}^Ty^{(N_H-1)}$ \eqref{R_i_T_effect}}
    }
    \ParFor(\Comment*[f]{$x^C\gets R_C x$}){$i=0, ..., N_{\mathcal{H}}-1$} {
      $x^{C}[i]\gets \Call{Sum}{x[k\bar l_{i} \slice k \bar l_{i + 1}-1]}$  \Comment*[r]{\eqref{R_C_effect}}
    }
    $y^{C}\gets \textsc{CoarseSolve}\!\left({x^{C}}\right)$  \Comment*[r]{$y^{C}\gets A_C^{-1}x^{C}$}
    $z\gets \Call{RepeatInterleave}{y^{C}, k\cdot l}$  \Comment*[r]{$z\gets R_C^T y^C$ \eqref{R_C_T_effect}}
    \Return $y + z$
  }
\end{algorithm}

\begin{algorithm}
  % When editing this algorithm, go through all references and check whether
  % the line numbers are correct.
  \caption{Hybrid Schwarz Preconditioner -- Multiplication By a Vector}
  \label{preconditioner-M-hyb-pseudocode}
  \KwReq{$k$, $l$, $\bar l$, \textsc{CoarseSolve}, \textsc{LocalSolve}$[0\slice N_H-1]$, $A_R$}

  \Function(\Comment*[f]{computes $M_{hyb}^{loc}x$, where $x \in \operatorname{range}(I - AR_C^TA_C^{-1}R_C)$}){M\_hyb\_loc($x$)} {
    \ParFor{$i=0, ..., N_H-1$} {
      $x^{(i)}\gets x[ki \slice k(i + 1)-1]$  \Comment*[r]{$x^{(i)}\gets R_ix$ \eqref{R_i_effect}}
      $y^{(i)}\gets \textsc{LocalSolve}[i]\!\left(x^{(i)}\right)$  \Comment*[r]{$y^{(i)}\gets A_i^{-1}x^{(i)}$}
      $y[ki \slice k(i + 1) - 1] \gets y^{(i)}$  \Comment*[r]{$y \gets R_0^Ty^{(0)}+ ...+R_{N_H-1}^Ty^{(N_H-1)}$ \eqref{R_i_T_effect}}
    }
    $z \gets \Call{MatVec}{A_R, y}$  \Comment*[r]{$z\gets R_CA y$}
    $w \gets \textsc{CoarseSolve}\!\left(z\right)$  \Comment*[r]{$w\gets A_C^{-1}z$}
    $t \gets \Call{RepeatInterleave}{w, k\cdot l}$  \Comment*[r]{$t\gets R_C^T w$ \eqref{R_C_T_effect}}
    \Return $y - t$
  }
\end{algorithm}
\section{Implementation}
\label{chap:implementation}

This section describes our GPU implementation of the domain decomposition
solver for the linear system~\eqref{auf}. We choose A100 SXM4 80GB as
the optimization target, and use it for all benchmarks and experiments.
However, the implementation supports any NVIDIA GPU with compute capability
8.0 or higher. Morever, we expect that the conclusions we draw from its
evaluation generalize to other highly parallel architectures.

The presented solver is implemented in Python 3 and relies on the PyTorch 2
framework~\cite{pytorch,pytorch2} for GPU-accelerated linear algebra operations.
This makes the implementation concise, free from memory management issues
and transportable across different hardware platforms supported by PyTorch.
Only when PyTorch operations prove insufficient, we implement custom CUDA kernels.
We use CUDA C++ \cite[Chapter 3]{cuda-guide} to
implement the batched dense matrix\nobreakdash-vector product
(see Section~\ref{sec:batched-dense-matrix-vector-product}),
since it requires careful handling of type casting and direct calls to
half-precision intrinsics. The majority of our kernels, however, are less low-level
and are therefore written in a subset of Python supported by the CUDA target
of a just-in-time compiler for Python -- Numba \cite{numba}.
This choice simplifies integration with PyTorch.
Moreover, static parameters of the kernels, such as data types or the BSR block size,
need not be specified in advance, but can be inferred from the input tensors.
Kernels are compiled on first invocation and then cached for all subsequent calls.

The complete code is included in the \texttt{dd\_solvers} package
available on GitHub [???]. It was tested on the NVIDIA Docker container tagged as
\texttt{nvcr.io/nvidia/pytorch:24.12-py3} \cite{pytorch-ngc}, which includes
Python 3.12.3, PyTorch 2.6.0 and NVIDIA CUDA 12.6.3. All additional solver
dependencies are specified in the accompanying \texttt{Dockerfile}.


\subsection{Solvers Overview}
\label{sec:solvers-overview}

\begin{figure}
  \centering
  \includesvg[width=\textwidth]{solvers-uml}
  \caption{Hierarchy of solvers. Solid lines indicate inheritance, dashed lines
  -- composition.}
  \label{fig:solvers-uml}
\end{figure}

The structure of the \texttt{dd\_solvers} package is illustrated in
Figure~\ref{fig:solvers-uml}. It defines three abstract base classes -- \verb|SparseSolver|,
\verb|DenseBatchSolver| and \verb|SparseBatchSolver| -- which specify
interfaces for solvers handling a single CSR matrix, a batch of dense matrices,
and a batch of CSR matrices, respectively. Because solvers of the same type share
a common interface, they can be interchanged easily, facilitating performance
comparisons.

At the core of the package is the abstract class \texttt{SchwarzOperator},
which encapsulates the shared components of the additive and hybrid Schwarz
preconditioners, particularly the setup phase (Algorithm~\ref{preconditioner-setup-pseudocode}).
The specific actions of $\bm{M_{add}}$ and $\bm{M_{hyb}^{loc}}$
(Algorithms~\ref{preconditioner-M-pseudocode} and~\ref{preconditioner-M-hyb-pseudocode})
are implemented in \texttt{AdditiveSchwarz} and \texttt{HybridSchwarz} classes, respectively;
both derive from \texttt{SchwarzOperator}. The details are discussed
in Section~\ref{sec:schwarz-preconditioners}.
Schwarz operators are parametrized by
the choice of coarse and local solvers: the coarse solver must subclass
\texttt{SparseSolver}, while the local solver
may be either a \texttt{DenseBatchSolver} or a \texttt{SparseBatchSolver}.
Dense batch solvers are preferable for small local problems (e.g.
$\mathcal{T}_H=\mathcal{T}_h$), whereas sparse solvers suit larger subdomains
(e.g. $\mathcal{T}_H=\mathcal{T}_{\mathcal{H}}$).

Since Schwarz operators are preconditioners, they must
be used in conjunction with an iterative method. Following Section~\ref{sec:additive-schwarz-method}
we employ the Preconditioned Conjugate Gradient (PCG) method.
As a native PCG implementation is available in neither PyTorch nor cuSPARSE, we
provide a custom implementation in the \texttt{CG} class.
Any \verb|SparseSolver| can serve as a~preconditioner; setting this parameter
to \verb|None| (default) yields an unpreconditioned CG solver.
The implementation is described in Section~\ref{sec:cg-method}.

In \verb|CUDSS| solver we provide a Python interface (with a C++ backend) to
the NVIDIA cuDSS library (v0.7.1\todo{sprawdz wersje})~\cite{cudss}, which
is a a high-performance direct sparse solver. As \verb|CUDSS| implements the \texttt{SparseSolver}
interface, it can be used either as the coarse solver within \texttt{SchwarzOperator}
or as a~standalone solver for \eqref{auf}.
NVIDIA cuDSS can employ more efficient algorithms when the matrix
is symmetric and positive definite. As it is the case for
the stiffness matrix~$\bm{A}$, as well as for the coarse solver matrix $\bm{A_C}$,
this option is enabled.

Available dense batch solvers for the \verb|SchwarzOperator| local subproblems are
\texttt{LU}, \texttt{Cholesky} and \texttt{Inv}, which implement LU decomposition,
Cholesky factorization and direct multiplication by the inverse, respectively.
The first two are thin wrappers around PyTorch's built-in batched solvers,
while \texttt{Inv} is our custom implementation of the explicit inverse
batch solver, described in Section~\ref{sec:explicit-inverse-batch-solver}.
The sole sparse batch solver, \texttt{BatchCUDSS}, leverages
the abilities of cuDSS to solve a batch of linear systems
with matrices in the CSR format. As in \texttt{CUDSS}, the matrices are assumed
to be symmetric and positive definite, which is the case for the local solvers'
matrices $\bm{A_0}, ..., \bm{A_{N_H-1}}$.

We also include a GPU-based algebraic multigrid preconditioner as our performance
baseline. This preconditioner wraps the NVIDIA AmgX library \cite{amgx}
via the \verb|pyamgx| bindings~\cite{pyamgx} and is provided by the \texttt{AMGX} class
in our package. Configuration details are discussed
in Section \ref{sec:amgx-preconditioner}.

As will be later demonstrated in Chapter~\ref{chap:experiments}, applying the preconditioner
in lower precision than the original matrix can deliver a significant performance
improvement. Accordingly, we supplement the Schwarz
preconditioners with two levels of precision control. First, the \texttt{precision}
attribute of the \texttt{SchwarzOperator} class may be set to either \verb|float64| or
\verb|float32|, which dictates the precision of the entire preconditioner
-- including both coarse and local solvers. Second, when the batch solver is a
direct inverse, the \verb|precision| attribute of the \texttt{Inv} class can
reduce the local solver's precision to 32 bits (\verb|float32|) or even 16 bits
(\verb|bfloat16| or \verb|float16|). We also allow the \verb|AMGX| preconditioner
to run in either single (\verb|float32|) or double (\verb|float64|) precision.

We conclude this section by illustrating the flexibility of the \texttt{dd\_solvers}
package with examples of constructing \texttt{SparseSolver} instances in
Listing~\ref{lst:creating-solvers}.

\begin{algorithm}
\vspace{-3pt}
\begin{lstlisting}[
  caption={Example \texttt{SparseSolver} instances.},
  label={lst:creating-solvers},
  style=Python
]
# Example 1: Direct solver using cuDSS.
s1 = CUDSS()

# Example 2: Unpreconditioned CG solver with default settings.
s2 = CG()

# Example 3: PCG with AmgX preconditioner applied in single precision.
s3 = CG(AMGX(precision=torch.float32))

# Example 4: PCG with hybrid Schwarz preconditioner using sparse
# direct solvers for both coarse and local problems.
s4 = CG(HybridSchwarz(
    coarse_solver=CUDSS(),
    local_solver=BatchCUDSS(),
))

# Example 5: PCG with additive Schwarz preconditioner applied in
# single precision, using dense local solver in a 16-bit precision.
s5 = CG(AdditiveSchwarz(
    precision=torch.float32,
    coarse_solver=CUDSS(),
    local_solver=Inv(precision=torch.float16),
))
\end{lstlisting}
\vspace{-3pt}
\end{algorithm}


\subsection{Schwarz Preconditioners}
\label{sec:schwarz-preconditioners}

The implementation of the additive and hybrid Schwarz preconditioners
follows Algorithms~\ref{preconditioner-setup-pseudocode}--\ref{preconditioner-M-hyb-pseudocode}.
We focus on the scenario where the solver is first initialized with
the system matrix $\bm{A}$ and the parameters $k$ and
$l=(l_0, ..., l_{N_{\mathcal{H}-1}})$, and then used to solve
a large number of systems with the same matrix $\bm{A}$ but different right-hand sides.
Consequently, we focus on describing the application of the preconditioners to
a vector. This operation is a part of the solution phase, and therefore constitutes
our optimization target. The implementation of the setup phase is a~straightforward
realization of Algorithm~\ref{preconditioner-setup-pseudocode} and is not
discussed further.

Thanks to the usage of high-level linear algebra operations provided by PyTorch,
the implementation of applying either preconditioner to a vector is concise enough
to be included here in its entirety. We present and discuss the implementation
of the additive Schwarz preconditioner (Listing~\ref{lst:preconditioner-solve}) in detail;
the hybrid Schwarz preconditioner is implemented similarly.

\begin{algorithm}
\vspace{-3pt}
\begin{lstlisting}[
  caption={Implementation of preconditioner's multiplication by a vector.},
  label={lst:preconditioner-solve},
  style=Python
]
@torch.compile
def solve(self, x: torch.Tensor) -> torch.Tensor:          # $x\mapsto M_{add}x$
  x_lower_precision = x.to(self.preconditioner_precision or x.dtype)
  x_i = x_lower_precision.reshape(self.n_solvers, -1)      # $x^{(i)} \gets R_i x$
  y_i = self.local_solvers.solve(x_i)                   # $y^{(i)} \gets A_i^{-1} x^{(i)}$
  y = y_i.flatten()                       # $\hspace{3pt} y \gets R_0^T y^{(0)} + ... + R_{N_H-1}^T y^{(N_H-1)}$
  x_solvers = x_lower_precision.reshape(self.n_solvers, -1).sum(dim=1)
  x_c = torch.segment_reduce(                             # $x^C \gets R_C x$
    x_solvers, reduce="sum", offsets=self.solvers_per_coarse_scan
  )
  y_c = self.coarse_solver.solve(x_c)                    # $y^C \gets A_C^{-1} x^C$
  y_solvers = y_c.repeat_interleave(
    self.solvers_per_coarse, output_size=self.n_solvers
  )
  y += y_solvers.repeat_interleave(                       # $y\ +\!= R_C^T y^C$
    self.dofs_per_solver, output_size=y.shape[0]
  )
  return y.to(x.dtype)
\end{lstlisting}
\vspace{-3pt}
\end{algorithm}

First, in line 3,
the $n$-element input vector $x$ is converted to lower precision if it is enabled
in the preconditioner configuration. The following lines 4--6 correspond to
lines 2--6 of Algorithm~\ref{preconditioner-M-pseudocode}. Since
vectors $x^{(i)}$ comprise blocks of $k$ consecutive elements of $x$, their
batched representation \verb|x_i| is obtained by simply \emph{reshaping}
the tensor $x$. In other words, \verb|x_i| shares memory with \verb|x_lower_precision|
and only indexes it differently; therefore no GPU computation is invoked in this step.
The local solvers are then executed in parallel to compute the batch of
vectors $y^{(i)}$, which is finally flattened to yield the intermediate result $y$.
Again, flattening does not involve any computation, as \verb|y| shares memory with \verb|y_i|.

The next step is to compute the vector $x^C=R_C x$ (cf.
Algorithm~\ref{preconditioner-M-pseudocode}, lines 7\nobreakdash--9).
Since this step sums consecutive blocks of
$kl_0, ..., kl_{N_H-1}$ elements of $x$, it can be split into two stages.
First, consecutive blocks of
$k$ elements of $x$ are summed (line 7) to form \verb|x_solvers| of
shape $(N_H,)$. Then, consecutive blocks of $l_0, ..., l_{N_H-1}$ elements of
\verb|x_solvers| are summed (lines 8--10) to produce \verb|x_c| of shape
$(N_{\mathcal{H}}, )$. The tensor \verb|solvers_per_coarse_scan| passed to the
\verb|segment_reduce| function should contain the cumulative sums
$\bar l_0, ..., \bar l_{N_{\mathcal{H}}-1}$ defined in \eqref{bar_l_i_def}.
Such two-stage reduction leverages the regularity of the step involving the
largest number of elements and was found to be more efficient than a single
reduction over the entire $x$. Note that this optimization relies on the
assumption of a constant number of unknowns per solver
(see Section~\ref{sec:preconditioner-alg}).

Then in line 11, the coarse solver is invoked to compute $y^C=A_C^{-1} x^C$.
Next, lines 12--17 compute $z=R_C^T y^C$ (cf. Algorithm~\ref{preconditioner-M-pseudocode}, line 15)
and add it to $y$ to yield the final result $y+z$.
Similar to the reduction of $x$ to $x^C$, this uses a two-step construction
that we found to be more efficient than a single scatter operation.
The result is finally converted back to the original precision in line 18.

The entire method is decorated with \verb|torch.compile|, which fuses compatible
operations into a~single kernel \cite{pytorch2}. For example, downcasting in line 3 and
summation in line 7 are combined to eliminate a redundant read of
\verb|x_lower_precision|\todo{czy to nadal jest prawda jak zmienilem kolejnosc?}.

Note that the preconditioners require a specific ordering of the unknowns,
as described in Section~\ref{sec:preconditioner-alg}. If this ordering cannot be
imposed during the discretization phase (which should be the preferred and
inexpensive approach for production use),
an additional preprocessing step is required to
permute the system matrix $\bm{A}$. Consequently, each solve also entails permuting
the right-hand side vector $\bm{f}$ prior to, and the solution vector $\bm{u^*}$
after a call to \verb|solve|. For such an occasion,
the \verb|dd_solvers| package supplies utility functions that perform these permutations
efficiently, given the three mappings: from fine elements to unknowns, from
fine elements to solver subdomains, and from solver subdomains to coarse elements.

\subsection{Explicit Inverse Batch Solver}
\label{sec:explicit-inverse-batch-solver}

Using the explicit inverse as a local solver in the Schwarz preconditioners
may appear counterintuitive, since it may have worse numerical properties and requires a more
expensive setup than standard solvers based on factorizations such as LU or Cholesky.
Nevertheless, multiplication by the inverse exposes more parallelism than
a triangular solve, as all elements of the output vector can be computed
independently, which is particularly beneficial for highly\nobreakdash-parallel architectures
such as GPUs. Therefore, we devoted additional effort to optimize the batch
inverse solver.

Its implementation is provided by the \verb|Inv| class. The
setup phase computes the batch of inverses
$\bm{A_0}^{-1}, ..., \bm{A_{N_H-1}}^{-1}$ with a single call to \verb|torch.linalg.inv|.
The local solvers' matrices and their inverses are stored as
three-dimensional tensors of shape $(N_H, k, k)$, where $k$ is the number of
unknowns per solver.

The solution phase reduces to computing the batched product (cf.
Algorithm~\ref{preconditioner-M-pseudocode}, line~4)
\[
  \bm{y^{(i)}}\gets \bm{A_i}^{-1}\bm{x^{(i)}},\quad i=0,...,N_H-1,
\]
where the batch of vectors $\bm{x^{(i)}}$ is provided as a two-dimensional tensor of shape
$(N_H, k)$. This can be easily achieved in PyTorch using the matrix multiplication
operator \verb|@| or the function \verb|torch.bmm|, both of which call
\verb|cublas<t>gemvStridedBatched| under the hood. Unfortunately, our
experiments revealed that this cuBLAS
function performs poorly for small matrix sizes ($k < 64$), which are typical for
the massively parallel variants of the ASM preconditioner. We therefore provide our own
batched dense matrix-vector product implementation, described below.

The \verb|Inv| class exposes a \verb|precision| attribute that permits internal
downcasting of the stored inverse matrices to lower precision, which can
improve the performance of the solution phase. Regardless of the internal
storage precision of the inverses, the output vectors $\bm{y^{(i)}}$ produced
by the solver always has the same precision as the input vectors $\bm{x^{(i)}}$ (as dictated by
the preconditioner), which may be either double (\verb|float64|) or
single (\verb|float32|).

\subsubsection{Batched Dense Matrix-Vector Product}
\label{sec:batched-dense-matrix-vector-product}

As an intermediate step in the \verb|Inv| implementation, we consider the following,
slightly more general problem:
given a batch of $N$ matrices of size $k \times k$ (stored as a tensor $mat$ of shape $(N, k, k)$)
and a batch of $N$ vectors of size $k$ (stored as a tensor $x$ of shape $(N, k)$),
compute the batched matrix-vector product and return the result as another tensor of
shape $(N, k)$.

As a~solution, we propose two algorithms implemented as CUDA kernels.
The first, denoted as RPT-I and presented in
Algorithm~\ref{alg:batched-mat-vec-irrespective}, implements a~straightforward approach.
Each thread computes a single element of the output tensor $y$ by alternately reading
elements from $mat$ and $x$ to accumulate the dot product. The second,
denoted as RPT-S, first loads all required elements of $x$ into shared memory
and then performs the dot product using these cached values. This approach is
detailed in Algorithm~\ref{alg:batched-mat-vec-synchronized}.

\begin{algorithm}
  \caption{Batched dense matrix-vector product -- RPT-I kernel}
  \label{alg:batched-mat-vec-irrespective}
  $t \gets blockDim.x \cdot blockIdx.x + threadIdx.x$\;
  \lIf {$t \geq N k$} { \Return }
  $b \gets \lfloor t / k \rfloor$\;
  $r \gets t \!\!\mod k$\;
  $acc \gets 0$\;
  \For{$c \gets 0;\; c < k;\; c\!+\!+$} {
    $mat\_elem \gets mat[b, c, r]$\;
    $mat\_elem\_higher\_prec \gets \Call{ConvertTo}{mat\_elem, \textbf{typeof}({x\_elem})}$\;
    $x\_elem \gets x[b, c]$\;
    $acc\ +\!= mat\_elem\_higher\_prec \cdot x\_elem$
  }
  $y[b, r] \gets acc$
\end{algorithm}

\begin{algorithm}
  \caption{Batched dense matrix-vector product -- RPT-S kernel}
  \label{alg:batched-mat-vec-synchronized}
  $batches\_per\_block \gets \lfloor blockDim.x / k \rfloor$\;
  $local\_batch\_idx \gets \lfloor threadIdx.x / k \rfloor$\;
  \lIf {$local\_batch\_idx \geq batches\_per\_block$} { \Return }

  $b \gets blockIdx.x \cdot batches\_per\_block + local\_batch\_idx$\;
  $r \gets threadIdx.x \!\!\mod k$\;
  \lIf {$b \geq N$} { \Return }

  $shared\_x[0 \slice blockDim.x - 1] \gets \text{init shared array of the same type as } x$\;
  $shared\_x[threadIdx.x] \gets x[b, r]$\;
  $sync\_threads()$\;

  $acc \gets 0$\;
  \For{$c \gets 0;\; c < k;\; c\!+\!+$} {
    $mat\_elem \gets mat[b, c, r]$\;
    $mat\_elem\_higher\_prec \gets \Call{ConvertTo}{mat\_elem, \textbf{typeof}({x\_elem})}$\;
    $x\_elem \gets shared\_x[local\_batch\_idx \cdot k + c]$\;
    $acc\ +\!= mat\_elem\_higher\_prec \cdot x\_elem$
  }
  $y[b, r] \gets acc$
\end{algorithm}

In RPT-I, the assignment of threads to output elements is independent of the batch
structure; thus any $k$ is supported regardless of the kernel block size.
We therefore use a constant block size of $128$ threads, which
was found to yield the best overall performance. In contrast, RPT-S requires
that each $k$-element vector from the batch is processed within a single thread block.
We set the block size to $256$ for $k\leq 128$, ensuring that at least two vectors
are processed per block. For $k>128$ we set the block size to smallest multiple of $32$
greater than or equal to $k$, such that each block processes a single vector.
This limits the maximum supported $k$ to $1024$; however, since such large matrices
are not used in Schwarz preconditioners, this limitation is acceptable.

The matrices $mat$ are assumed to be
stored in column-major order (cf. lines 7 and 12 in
Algorithms~\ref{alg:batched-mat-vec-irrespective} and~\ref{alg:batched-mat-vec-synchronized},
respectively), which ensures that most accesses to the
elements of $mat$ can be coalesced (see \cite[Section 5.3.2]{cuda-guide}). This
matches the \verb|Inv| solver use case, since the matrices $\bm{A_i}^{-1}$
returned by \verb|torch.linalg.inv| are already in column-major order.
Even if they were not, no additional steps would be required, because the matrices $\bm{A_i}$
(and therefore their inverses $\bm{A_i}^{-1}$) are symmetric (see~\eqref{A_i_def}).

The matrices $mat$ may be stored in a precision lower than that of vectors $x$ and $y$.
In such cases each matrix element is upcast after being
fetched from memory. Thus multiplication
and accumulation are always performed in the precision of vectors $x$ and $y$.
Because these are only two floating\nobreakdash-point operations per element read from $mat$,
the algorithm is memory-bound and reducing computation precision does not
improve performance, which our experiments confirm.

As is usual for memory-bound operations, we measure the kernel performance
in terms of achieved memory throughput, calculated
as the ratio of the total size of the input and output tensors ($mat$, $x$, $y$)
to the kernel execution time.
To assess the optimality of our algorithms, we compare the achieved throughput
to the peak memory bandwidth measured on the target GPU using the
\verb|triad| benchmark from the BabelStream suite \cite{babelstream}.

We benchmarked our algorithms for various combinations of precisions.
For $mat$ we used $\verb|float16|$, $\verb|bfloat16|$, $\verb|float32|$
and $\verb|float64|$ when $x$ was in $\verb|float64|$, and
$\verb|float16|$, $\verb|bfloat16|$, $\verb|float32|$ when $x$ was in $\verb|float32|$.
We tested matrix sizes $k$ ranging from $2$ to $149$, which covers the
typical sizes used in the Schwarz preconditioners (see
Section~\ref{chap:experiments}).
For each $k$ we selected
the largest batch size $N$ satisfying $Nk^2\leq 2^{27}$
to saturate the GPU. The PyTorch implementation serves as a~baseline.
Note that PyTorch requires $mat$ and $x$ to have the same precision;
in mixed\nobreakdash-precision cases we therefore need to downcast $x$ before
multiplication and upcast the result $y$ afterwards. This introduces extra
overhead and yields much less accurate outputs than our implementation.

The results are similar for all precision combinasions; therefore, in
Figure~\ref{fig:matrix_vec_tput} we present only
two representative cases: when $mat$ is stored in \verb|float16| and
$x$ is in either \verb|float32| or \verb|float64|.
RPT-S sustains high throughput across the tested matrix sizes, reaching
over $84\%$ of the achievable A100 memory bandwidth.
RPT-I outperforms RPT-S only for a few smallest matrix sizes, where the
overhead of synchronizing threads in RPT-S becomes significant.
For $k\geq 16$, RPT-S is consistently faster
than RPT-I, and for some precision combinations (for example
float16$\times$float64) the difference is substantial.

\begin{figure}
  \centering
  \includesvg[width=\textwidth]{matrix_vec_tput}
  \caption{Performance of the batched dense matrix-vector product at different
  precisions: our implementation vs. PyTorch. The red line indicates the
  measured peak memory bandwidth of the A100 SXM4 80GB used in the experiments.}
  \label{fig:matrix_vec_tput}
\end{figure}

In contrast, the
PyTorch implementation performs poorly at the smallest matrix sizes, being
$10-70\times$ slower than ours for $k\leq 4$. Its throughput grows with $k$
and becomes competitive around $k\approx 64$, with the exact threshold
depending on the precision. However, the Schwarz preconditioners primarily targets
smaller matrix sizes when a~\verb|DenseBatchSolver| is used. In
Section~\ref{chap:experiments} we show that, in most scenarios, the optimal
choice is $21\leq k\leq 60$\todo{czy na pewno?} with \verb|float16| precision
-- settings for which RPT-S is $1.43 - 4.75\times$ faster than PyTorch.
Therefore we employ our kernels in the \verb|Inv| class, choosing RPT-I
for smallest matrices ($k\leq 5$ or $k\leq 15$ depending on the precision)
and RPT-S otherwise.

\subsection{BSR Matrix-Vector Product}

The per-iteration cost of the PCG method is dominated by two operations:
the preconditioner solve and the sparse matrix-vector product. The latter can become
the most expensive step when using a lightweight preconditioner (such as massively
parallel variant of the Schwarz method), as confirmed by our benchmarks.

We therefore optimize the computation of the product with the stiffness matrix $\bm{A}$
by exploiting its structure. If the unknowns associated with each element
form a sequence of consecutive entries (which is assured by the proper permutation), then the
matrix $\bm{A}$ has a block-sparse structure, with block size $b$ equal to the number
of unknowns per element. Specifically, the nonzero entries of $\bm{A}$ can be grouped
into dense square blocks of size $b\times b$, where the nonzero blocks themselves form
a sparse pattern. Consequently, $\bm{A}$ can be represented in the Block Sparse Row (BSR)
format \cite[Section 3.3.6]{cusparse}. Compared to CSR, the BSR format reduces the size
of the row-offset array by a factor of $b$ and the column-index array by a factor of $b^2$.

Beyond memory savings, BSR format accelerates the matrix-vector
product. To leverage this, the \verb|CG| class has an extra
attribute \verb|bsr_matmul|, which enables the conversion of the matrix $\bm{A}$
to the BSR format during the setup phase.

Multiple BSR matrix-vector algorithms are available.
The simplest approach is to use \texttt{cusparse<t>bsrmv} routine from the
cuSPARSE library, which is accessible via a PyTorch
interface. However, in our benchmarks it underperformed the CSR matrix-vector
product from the same library.
To improve performance, we implemented (in CUDA kernels)
four algorithms from~\cite{bsr-mv}: block-by-block,
column-by-column, row\nobreakdash-per\nobreakdash-thread without synchronization (RPT) and
row-per-thread with synchronization (RPT-S). Reductions in the block-by-block and
column\nobreakdash-by\nobreakdash-column kernels use warp-shuffle operations
\cite[Section 7.22]{cuda-guide}. All implemented kernels
run on a one-dimensional grid, with block size 32 for the RPT and RPT-S
kernels, and an experimentally chosen 256 for the block-by-block and column\nobreakdash-by\nobreakdash-column
kernels.

We also introduce a modified RPT variant: RPT-I, which assigns
rows to threads irrespective of thread block boundaries (see
Algorithm~\ref{alg:bsr-mv-naive}). In the pseudocode, the BSR matrix $A$ is
represented by the arrays $A.crow\_indices$, $A.col\_indices$ and $A.values$.
The blocks of $A.values$ are assumed to be stored in column-major order
within each block (cf. line 10), which follows~\cite{bsr-mv} and ensures
coalesced memory accesses. The algorithm computes the matrix-vector product
$Ax$ and stores the result in the $global\_out$ array. The RPT-I kernel likewise
executes on a one-dimensional grid and uses a~block size of 256 threads.

\begin{algorithm}
  \caption{BSR matrix-vector product -- RPT-I kernel}
  \label{alg:bsr-mv-naive}
  $target\_row \gets blockDim.x \cdot blockIdx.x + threadIdx.x$\;
  $target\_block\_row \gets \lfloor target\_row\ /\ b\rfloor$\;
  \lIf {$target\_block\_row \geq A.crow\_indices.size - 1$} { \Return }
  $r \gets target\_row\!\!\mod b$\;
  $first\_block \gets A.crow\_indices[target\_block\_row]$\;
  $after\_last\_block \gets A.crow\_indices[target\_block\_row + 1]$\;
  $local\_out \gets 0$\;
  \For{$block \gets first\_block;\; block < after\_last\_block;\; block\!+\!+$} {
    \For{$c\gets 0;\; c < b;\; c\!+\!+$} {
        $A\_elem \gets A.values[block][c][r]$\;
        $x\_elem \gets x[A.col\_indices[block] \cdot bs + c]$\;
        $local\_out\ +\!= A\_elem \cdot x\_elem$
    }
  }
  $global\_out[target\_row] \gets local\_out$
\end{algorithm}

All algorithms were benchmarked using stiffness matrices produced
by the Discontinuous Galerkin discretization described in Section~\ref{sec:discretization}.
We tested polynomial degrees $p=1,...,5$ in 2D and $p=1,2,3$ in 3D, yielding
block sizes ranging from $3$ to $21$. Indices were stored as 32-bit integers,
while matrix entries and all computations used 64-bit floating-point precision.
Throughput was measured on an NVIDIA A100 GPU and, as before, compared to the peak
memory bandwidth measured with
the \verb|triad| benchmark from the BabelStream suite \cite{babelstream}.
The results are shown in in Figure~\ref{fig:bsrmv_memory_throughput}.

\begin{figure}
  \centering
  \includesvg[width=\textwidth]{bsrmv_memory_throughput}
  \caption{Performance of our implementations of the BSR matrix-vector
  product for various discretization parameters. Block-by-block algorithm requires
  $b\leq 5$.}
  \label{fig:bsrmv_memory_throughput}
\end{figure}

The \texttt{cusparse<t>bsrmv} routine, which we use as a baseline, was competetive
only for $b=4$ and $b=15$. In the remaining cases it performed poorly, being as much as
$2.68\times$ slower than our implementations.
We suspect that the cuSPARSE algorithm is optimized for block sizes being
close to powers of two, and is therefore not well suited for our use case.

Relative performance of the algorithms
introduced in \cite{bsr-mv} (block-by-block, column-by\nobreakdash-column, RPT and RPT-S),
agrees with the conclusions reported there. Absolute throughput numbers
are substantially higher than in \cite{bsr-mv}, as expected given the Kepler
hardware used in that study versus much more advanced Ampere architecture
of the A100 GPU employed here. Nevertheless,
on the A100 all these algorithms are outperformed by the straightforward
RPT\nobreakdash-I approach, which is the fastest method for all tested scenarios.
In addition, it attains very high memory throughput, ranging from $77\%$ to $92\%$
of the measured A100 peak bandwidth. Combined with the memory savings
of the BSR format, this yields a $1.42-1.69\times$ speedup over the
CSR matrix-vector product from cuSPARSE. For these reasons, we use the BSR format
and the RPT-I algorithm as the default configuration in our PCG implementation.

\subsection{Reference AmgX Preconditioner}
\label{sec:amgx-preconditioner}

AmgX \cite{amgx, amgx-reference} is a high-performance NVIDIA library for solving
large sparse linear systems on NVIDIA GPUs. It provides a variety of
algebraic multigrid (AMG) preconditioners and Krylov subspace solvers which
can be flexibly combined and configured. Due to the large number of possible configurations,
careful selection is necessary to achieve optimal performance for a given problem
and, in our case, to ensure a "fair" comparison with our ASM implementation.

We followed the procedure below to select an AmgX configuration. Starting from all
example configurations provided with AmgX (\verb|AMGX/src/configs|), we modified any that used
GMRES or FGMRES to use PCG instead, and employed a stopping criterion that reduces
the initial Euclidean norm of the residual by a factor of $10^9$. Benchmarks
used system \eqref{auf} obtained for $p=1$ and $\mathcal{T}_h$ being a uniform
simplicial mesh ($\mathcal{S}_m^d$) of a unit square ($d=2$) or a~unit cube ($d=3$);
see Section~\ref{sec:experimental-setup} for details on the meshes.

We first tested small problems: 12k DoFs in 2D and 25k DoFs in 3D
($\mathcal{T}_h\in \{\mathcal{S}_6^2, \mathcal{S}_3^3\}$), to identify
configurations that successfully converge. We then benchmarked the convergent configurations
on medium problems: 393k DoFs in 2D and 197k DoFs in 3D
($\mathcal{T}_h\in \{\mathcal{S}_8^2, \mathcal{S}_5^3\}$), and retained those
that were at most 50\% slower than the fastest one in each dimension.
Finally, we ran the selected configurations on large problems:
25M DoFs in 2D and 50M DoFs in 3D ($\mathcal{T}_h\in \{\mathcal{S}_{11}^2, \mathcal{S}_7^3\}$).
The configuration named \verb|GMRES_AMG_D2| and all configurations that
used the HMIS coarsening selector were excluded as they produced runtime errors
on the large 3D problem.  This left four configurations (the same set in both dimensions),
which we list in Tables~\ref{tab:2d-amgx-configs} and~\ref{tab:3d-amgx-configs}.

% \begin{table}
%   \centering
%   \caption{Performance of AmgX configurations on A100: 2D problem with 25M DOFs.}
%   \label{tab:2d-amgx-configs}
%   \input{tables/amgx_2d.tex}
% \end{table}

% \begin{table}
%   \centering
%   \caption{Performance of AmgX configurations on A100: 3D problem with 50M DOFs.}
%   \label{tab:3d-amgx-configs}
%   \input{tables/amgx_3d.tex}
% \end{table}

All selected configurations employ a classical algebraic multigrid preconditioner
with a~V\nobreakdash-cycle and parallel coarsening algorithm (PMIS). They differ
in the choice of smoother and interpolation operator. See the AmgX
documentation \cite{amgx-reference} for details on these options and their
parameters. \verb|AMG_CLASSICAL_AGGRESSIVE_L1| is the fastest configuration in both
dimensions, but we retain all four for the experiments in
Chapter~\ref{chap:experiments}, since those scenarios include larger $p$
and mixed-precision arithmetic, which may affect the relative performance of
the configurations.
Note that we did not further tune AmgX settings here; additional performance
gains may be possible through parameter optimization.

The preconditioners used in the selected configurations are exposed in the
\verb|dd_solvers| module via the \verb|AmgX| class. Its constructor accepts
a configuration name, which might be one of
\verb|AGGRESSIVE_L1|, \verb|L1_TRUNC|, \verb|AGGRESSIVE_L1_TRUNC|, or
\verb|AGGRESSIVE_CHEB_L1_TRUNC| (we omit the common \verb|AMG_CLASSICAL_| prefix
for brevity). Using \verb|AMGX| together with our \verb|CG| class
corresponds to the original AmgX configurations listed in
Tables~\ref{tab:2d-amgx-configs} and~\ref{tab:3d-amgx-configs}, except that the
AmgX PCG solver is replaced by our PCG implementation.

\section{Parameter Tuning and Performance Evaluation}
\label{chap:experiments}

In this chapter we experimentally evaluate the performance of our implementation
of the ASM solver described in Chapter~\ref{chap:implementation}. The experimental
setup is detailed in Section~\ref{sec:experimental-setup}. We first verify
theoretical properties of the preconditioned ASM operator in
Section~\ref{sec:dependence-on-parameters}. Next, Section~\ref{sec:precision-experiments}
examines the effect of using lower precision arithmetic,
and Section~\ref{sec:experiments-parallelism} investigates the impact of
of the coarse and local solvers' mesh choices. Finally,
Section~\ref{sec:experiments-amgx-comparison} compares our implementation to
the reference AmgX preconditioner.

\subsection{Experimental Setup}
\label{sec:experimental-setup}

In all experiments we consider $\Omega=(0, 1)^d$ with $d\in \{2, 3\}$ and fix
$\rho(x)=1+\sum_{i=1}^d x_i^2$. We set $\delta=7$ for $d=2$ and $\delta=17$ for $d=3$.
The right-hand side vector $\bm{f}$ is manufactured such that
$u(x)=\prod_{i=1}^d \sin(4\pi x_i)$ is the exact solution of \eqref{div_eq}.
Note that the experiments are conducted in the general three-grid scenario
with $p\geq 1$, which is not covered by \cite{Przykry} nor~\cite{asm-hp-two-meshes},
but was partially addressed in~\cite{Przykry-hp}.
In this setting we expect the condition number of the preconditioned operator
$T_{add}$ to satisfy
\[
  \kappa(T_{add})\lesssim p^2\frac{\mathcal{H}^2}{hH},
\]
where by $\lesssim$ we mean inequality up to a constant independent of
$h$, $H$, $\mathcal{H}$, $p$ and jumps of the coefficient $\rho$. Moreover,
from the abstract theory of Schwarz operators (see~\cite{toselli}) it follows that
\[
  \kappa(T_{hyb})\lesssim \kappa(T_{add}).
\]

We\todo{uwaga o tym, ze wystarczy ograniczyc sie do takich siatek i ciaglego wsp}
introduce two types of uniform meshes for use in the experiments: consisting
of either only hypercubes or only simplices. For $d\in\{2, 3\}$ and integer $m\geq 0$, let
$\mathcal{C}_m^d$ denote the partition of $(0, 1)^d$ into $2^{md}$ equal hypercubes
(squares in 2D, cubes in 3D) of side length~$2^{-m}$. The mesh $\mathcal{S}_m^d$
is obtained by subdividing each hypercube of $\mathcal{C}_m^d$ into simplices:
by the natural diagonal subdivision into $2$ triangles in 2D, and by the
Coxeter-Freudenthal-Kuhn triangulation (see, e.g., \cite{tetrahedral-subdivision})
into $6$ tetrahedra in 3D. When the dimension is clear from the context we
drop the superscript and write $\mathcal{C}_m$ and $\mathcal{S}_m$.

Examples for $d=2$ are shown in Figure~\ref{fig:mesh_examples}. In $\mathcal{S}_m^d$
(for both $d=2,3$) we use an alternating diagonal orientation across neighboring
cells. Our experiments have shown that this alternating subdivision yields
smaller approximation errors and faster convergence of
all considered solvers than the non-alternating one. The alternation
is chosen consistently for all~$m$, so that
$\mathcal{S}_m^d\subseteq \mathcal{S}_{m+1}^d$ for $m\geq 0$. Moreover
$\mathcal{C}_m^d \subseteq \mathcal{C}_{m + 1}^d$ and $\mathcal{C}_m^d \subseteq \mathcal{S}_m^d$
for all $m\geq 0$.

We always use $\mathcal{S}_m$ (for some $m\geq 0$) as the fine mesh $\mathcal{T}_h$, which gives
$h=2^{-m}\sqrt{d}$. For any integers $0\leq \mathcal{M}\leq M \leq m$ we select
the coarse and local solvers' meshes from the following combinations:
\begin{itemize}
  \item $\mathcal{T}_\mathcal{H}=\mathcal{S}_\mathcal{M}$, $\mathcal{T}_H=\mathcal{S}_M$;
  \item $\mathcal{T}_\mathcal{H}=\mathcal{C}_\mathcal{M}$, $\mathcal{T}_H=\mathcal{S}_M$;
  \item $\mathcal{T}_\mathcal{H}=\mathcal{C}_\mathcal{M}$, $\mathcal{T}_H=\mathcal{C}_M$.
\end{itemize}
In all cases $\mathcal{T}_\mathcal{H}\subseteq \mathcal{T}_H\subseteq \mathcal{T}_h$ with
$\mathcal{H}=2^{-\mathcal{M}}\sqrt{d}$ and $H=2^{-M}\sqrt{d}$. This choice
of $\mathcal{T}_\mathcal{H}$ and $\mathcal{T}_H$ satisfies the regularity conditions
of Theorems~\ref{thm:Przykry-asm} and~\ref{thm:two-meshes-asm}, as well as the
assumption of an equal number of unknowns per local solver
(see Section~\ref{sec:preconditioner-alg}).

The code for the experiments is included in the \verb|experiments| package
attached to this thesis.
Both mesh generation and discretization are performed on a CPU using DOLFINx 0.9.0~\cite{dolfinx}.
The resulting stiffness matrix $\bm{A}$ and load vector $\bm{f}$ from the discrete system
\eqref{auf} are transferred to the GPU, which performs both the setup
and the solution phase of the chosen solver. The computed solution of \eqref{auf}
may finally be transferred back to the CPU for further processing. The resulting solution
is always computed in double precision, but lower\nobreakdash-precision arithmetic
might be used internally by the preconditioner according to the choice of
its parameters (see Section~\ref{sec:solvers-overview}).

The GPU we use to run all experiments in this chapter is an NVIDIA A100 SXM4 80GB.
We use CUDA Events \cite[Section 3.2.8.8]{cuda-guide} to measure GPU timings.
The general measurement procedure is as follows. We perform $N$ rounds, each
consisting of $K$ consecutive runs of the measured phase and record the total
time of each round independently. For each round we compute the per-run time
as the round time divided by $K$, and we report the minimum per-run time over
all rounds.

Executing multiple rounds is necessary for several reasons. First, after
idle periods the GPU requires a warm-up interval to reach its maximum frequency.
Second, the initial run may be significantly slower due to just-in-time
kernel compilation. Finally, individual runs can be slower because of external
factors such as other users launching tasks on the same host. On the other hand,
choosing $K>1$ is useful only when the measured function is very short
(below roughly $10$\si{ms}), so that the measurement overhead is non-negligible
and must be amortized by repeated execution. Accordingly, we use the following
settings: $N=3$, $K=1$ for setup phases, $N=30$, $K=1$ for solve-time measurements,
and $N=30$, $K=10$ when measuring a single execution of the ASM preconditioner
or one of its constituent parts.

\begin{figure}
  \centering
  \includesvg[width=0.7\textwidth]{mesh_examples}
  \caption{Example experimental meshes and their refinement relationships for $d=2$:
  ${\mathcal{S}_1\subseteq \mathcal{S}_2\subseteq \mathcal{S}_3}$ (left) and
  $\mathcal{C}_1\subseteq \mathcal{C}_3\subseteq \mathcal{S}_3$ (right).}
  \label{fig:mesh_examples}
\end{figure}

\subsection{Exploiting Mixed Precision}
\label{sec:precision-experiments}

In this section we study the effectiveness of solving the linear system \eqref{auf}
using \emph{mixed precision}: PCG iteration is executed in double precision
to preserve accuracy, while the ASM preconditioner uses lower precisions internally
to accelerate the computation. We test this approach on three discrete systems,
obtained for $d=2$ and $(p, \mathcal{T}_h)\in \{(1, \mathcal{S}_{10}), (3, \mathcal{S}_9),
(5, \mathcal{S}_8)\}$. For each case we take
$\mathcal{T}_H=\mathcal{T}_{\mathcal{H}}=\mathcal{S}_\mathcal{M}$ with
$m-2\leq \mathcal{M}\leq m$.

We use \verb|Inv| as the local solver, as it offers most flexibility for precision
control. There are four possible configurations
(Listing~\ref{lst:precision-experiments-solvers}): the entire ASM preconditioner
in double precision (\verb|fp64|), in single precision (\verb|fp32|)
and in single precision with the local solvers using 16-bit types
(\verb|fp32+fp16| or \verb|fp32+bf16|).

\begin{algorithm}
\vspace{-3pt}
\begin{lstlisting}[
  caption={Solvers used in the mixed-precision experiments.},
  label={lst:precision-experiments-solvers},
  style=Python
]
CG(preconditioner, rtol=1e-9)
for preconditioner in (
  ASM(torch.float64, Inv(torch.float32), CUDSS()),   # fp64
  ASM(torch.float32, Inv(torch.float32), CUDSS()),   # fp32
  ASM(torch.float32, Inv(torch.float16), CUDSS()),   # fp32+fp16
  ASM(torch.float32, Inv(torch.bfloat16), CUDSS()),  # fp32+bf16
)
\end{lstlisting}
\vspace{-3pt}
\end{algorithm}

PCG iteration counts required to reduce the initial Euclidean norm of the
residual by a~factor of $10^9$ are reported in
Table~\ref{tab:experiment-precision-iters}. Using \verb|fp32|
did not change the number of iterations except in one case ($p=5$, $m=8$, $M=\mathcal{M}=6$).
The \verb|fp32+fp16| configuration added at most a few iterations for $p=3$ and $p=5$.
We attribute such good behaviour of the \verb|fp32+fp16| variant to the fact that
half precision is used only to store the inverses of the local solvers' matrices,
while their application to batches of vectors is performed in single precision
(see Section~\ref{sec:explicit-inverse-batch-solver}).

By contrast, using \verb|fp32+bf16| significantly degrades convergence, requiring
up to $2.5\times$ more iterations than \verb|fp64| in the extreme case $p=5$, $m=8$, $M=\mathcal{M}=6$.
The comparison of the convergence behaviour for this scenario is shown in
Figure~\ref{fig:precision-residuals}. While \verb|fp64|, \verb|fp32|
and \verb|fp32+fp16| converge at the same steady rate, the \verb|fp32+bf16| run begins
to slow when the residual norm reaches approximately $10^{-5}$ and eventually
plateaus near $10^{-9}$.

\begin{figure}
  \centering
  \includesvg[width=\textwidth]{precision_residuals}
  \caption{Euclidean norm of the residual at each PCG iteration for different
  preconditioner precision settings (selected case). Parameters:
  $d=2$, $p=5$, $\mathcal{T}_h=\mathcal{S}_8$ and
  $\mathcal{T}_H=\mathcal{T}_\mathcal{H}=\mathcal{S}_6$.}
  \label{fig:precision-residuals}
\end{figure}

These findings are consistent with the timings reported in Figure~\ref{fig:precision-plots}.
We measured not only the total solution time but also estimated the time spent
in each ASM component by timing them independently and scaling by the number of
PCG iterations. Local solvers benefit substantially from lower precision,
with speedups of $1.75 - 1.96\times$ for \verb|fp32| and $2.62 - 3.40\times$ for
\verb|fp32+fp16|, consistent with Section~\ref{sec:batched-dense-matrix-vector-product}.
On the other hand, the time spent in the coarse \verb|CUDSS| solver is only marginally
affected, showing speedups of $1.09 - 1.20\times$ when using \verb|float32| instead
of \verb|float64|. This is expected since \verb|CUDSS| employs a sparse
decomposition algorithm that relies on reading the integer indices of the
nonzero entries, which are unaffected by floating-point precision changes.

\begin{figure}
  \centering
  \includesvg[width=\textwidth]{precision_plots}
  \caption{Comparison of the solution time, decomposed by component, for various
  problems and preconditioner precision settings. Parameters:
  $d=2$, $\mathcal{T}_h=\mathcal{S}_m$ and
  $\mathcal{T}_H=\mathcal{T}_\mathcal{H}=\mathcal{S}_\mathcal{M}$.}
  \label{fig:precision-plots}
\end{figure}

The remainder of the ASM time -- application of the restriction and prolongation
matrices $\bm{R_C}$ and $\bm{R_C}^T$, and summation of the solvers' contributions
-- is also insensitive to the precision settings. There are two possible reasons
for this. First, restriction and prolongation to the coarse space $V_C$ require
reading integer index sequences $l_i$ and $\bar l_i$ that determine the nonzero
structure of $\bm{R_C}$ and $\bm{R_C}^T$ (see Algorithm~\ref{preconditioner-M-pseudocode});
these indices are not subject to precision reduction. Second, this part incurs
type conversions when lower precisions are used, which reduce the attainable
speedup.

Finally, time spent outside the ASM preconditioner is essentially unaffected,
as expected, since the remainder of the PCG iteration is always performed in
double precision, and the modest increases in iteration counts have a negligible
impact on this cost. The only exception is the \verb|fp32+bf16| case for $p=5$,
discussed above.

From these observations we conclude that the largest performance
gains from using lower precisions may be achieved when the local solvers dominate
the preconditioner cost, i.e., when the local problems are large and the coarse
problem is small. Indeed, Figure~\ref{fig:precision-plots} shows that, for smaller
$M=\mathcal{M}$ and larger $p$, the speedup in the total solution time obtained
with \verb|fp32| and \verb|fp32+fp16| over \verb|fp64| increases.

Overall, our measurements indicate that mixed precision is highly beneficial.
With \verb|Inv| as the local solver, we observed speedups up to $1.44\times$
for \verb|fp32| and $1.81\times$ for \verb|fp32+fp16|. The \verb|fp32+fp16| variant
should, however, be used with caution. When the local problems are small, the
additional reduction in local solvers' time offered by \verb|fp32+fp16| relative
to \verb|fp32| may be insufficient to introduce a meaningful speedup in total solution time
(this was observed for $\mathcal{M}=M=m$). Even though this was not observed here,
\verb|fp32+fp16| may, in extreme cases, also cause a minor slowdown compared
with \verb|fp32|, as we will see in Section~\ref{sec:experiments-parallelism}.
The \verb|fp32+bf16| variant should be avoided, as it offers no performance
benefit over \verb|fp32+fp16| and may severely degrade convergence.

\subsection{Tuning Parallelism by Selecting $\mathcal{T}_\mathcal{H}$ and $\mathcal{T}_H$}
\label{sec:experiments-parallelism}

This section examines how the performance of the ASM solver depends on the choice
of the coarse mesh $\mathcal{T}_\mathcal{H}$ and the local solvers' mesh $\mathcal{T}_H$.
Our metric of interest is now total solution time rather than the
iteration count considered in Section~\ref{sec:dependence-on-parameters}.

We set $d=2$ and $p=1$ to maximize the amount of possible parallelism: the local solvers' matrices
can be as small as $3\times 3$ when $\mathcal{T}_H=\mathcal{T}_h$. We fix
$\mathcal{T}_h=\mathcal{S}_{11}$ as the fine mesh, which yields a discrete problem
with roughly 25 million degrees of freedom. For the coarse mesh we consider
$\mathcal{C}_\mathcal{M}$ and $\mathcal{S}_\mathcal{M}$ with $8\leq \mathcal{M}\leq 11$.
For each choice of $\mathcal{T}_\mathcal{H}$ we evaluate all $\mathcal{T}_H$
satisfying $\mathcal{T}_\mathcal{H}\subseteq \mathcal{T}_H\subseteq \mathcal{T}_h$,
and being either $\mathcal{C}_M$ or $\mathcal{S}_M$ for some $M$.

We test all available local solvers: \verb|Inv|, \verb|LU|, \verb|Cholesky| and
\verb|BatchCUDSS|. In all cases we apply the preconditioner in single precision,
as the previous section demonstrated its advantage over double precision.
For \verb|Inv| we additionally evaluate the half-precision configuration
\verb|fp32+fp16|. All timings report the time required to reduce the initial
Euclidean norm of the residual by a~factor of $10^9$. Precise solver configurations
are given in Listing~\ref{lst:solver-experiments-solvers}.

\begin{algorithm}
\vspace{-3pt}
\begin{lstlisting}[
  caption={Solvers used in the parallelism experiments.},
  label={lst:solver-experiments-solvers},
  style=Python
]
CG(ASM(torch.float32, local_solver, CUDSS()), rtol=1e-9)
for local_solver in (
  Inv(),
  Inv(torch.float16),
  LU(),
  Cholesky(),
  BatchCUDSS(),
)
\end{lstlisting}
\vspace{-3pt}
\end{algorithm}

The comparison of solution times is shown in Figure~\ref{fig:solver-plots}.
Following the approach from the previous section, we also estimated the time
contribution of each ASM component. Missing entries correspond to
a case when either some dense batch solvers failed because local
problems were too large ($\mathcal{T}_H\in \{\mathcal{C}_8, \mathcal{S}_8\}$),
or the sparse batch solver failed due to the excessive number of local problems
($\mathcal{T}_H \in \{\mathcal{C}_{11}, \mathcal{S}_{11}\}$).
We adopt \verb|Inv(float16)| as the reference local solver since it was the
top variant in nearly all cases; exceptions occured for
$\mathcal{T}_H\in \{\mathcal{C}_{11}, \mathcal{S}_{11}\}$, where single-precision
\verb|Inv()| was marginally faster.

% \begin{figure}
%   \centering
%   \includesvg[width=\textwidth]{solver_plots}
%   \caption{Comparison of the \texttt{Inv(float16)} solution time, decomposed by component,
%   for various $\mathcal{T}_\mathcal{H}$ and $\mathcal{T}_H$. Markers indicate
%   the solution time when \texttt{Inv(float16)} is replaced by other local solvers.
%   Parameters: $d=2$, $p=1$ and $\mathcal{T}_h=\mathcal{S}_{11}$.}
%   \label{fig:solver-plots}
% \end{figure}

The other dense batch solvers (\verb|LU| and \verb|Cholesky|) reported in
Figure~\ref{fig:solver-plots} are noticeably
slower than the direct inverse (note the logarithmic scale for times above 1\,s).
This is expected since they are harder to parallelize effectively, moreover,
we invested additional effort in optimizing \verb|Inv| relative to the baseline
PyTorch implementation (see Section~\ref{sec:explicit-inverse-batch-solver}).
Although \verb|LU| and \verb|Cholesky| offer faster setup times
than \verb|Inv| for large local problems (Figure~\ref{fig:local-solver-setup-times}),
their slower solve phase outweighs that benefit, even when only a single
solve is performed per setup.

% \begin{figure}
%   \centering
%   \includesvg[width=\textwidth]{local_solvers_setup_times}
%   \caption{Setup time of the local solvers (including the construction of
%   $\bm{A_0}, ..., \bm{A_{N_H-1}}$ from~$\bm{A}$) as a~function of $\mathcal{T}_H$
%   for the dense batch solvers (left) and the sparse batch solver (right). Parameters:
%   $d=2$, $p=1$ and $\mathcal{T}_h=\mathcal{S}_{11}$.}
%   \label{fig:local-solver-setup-times}
% \end{figure}

The \verb|BatchCUDSS| solver performs poorly for small local problems,
which is expected. Surprisingly, it only marginally outperforms \verb|LU()| even
for $T_H=\mathcal{C}_8$, where the local solvers' matrices are $384\times 384$.
Moreover, the setup times for \verb|BatchCUDSS| are two orders of magnitude larger
than those of the dense batch solvers. \verb|BatchCUDSS| would definitely be
preferable for larger local problems, when the dense batch solvers would exhaust
GPU memory. However, we find it very unlikely to provide competitive solution
time in these cases, given that its best performance among considered scenarios
occured for $\mathcal{T}_\mathcal{H}=\mathcal{T}_H=\mathcal{C}_{10}$, and
worsened as $\mathcal{T}_\mathcal{H}$ became coarser.

The above observations indicate that the explicit inverse dense batch solver
is the optimal choice for problems of the size considered here. We now compare
its performance across different choices of $\mathcal{T}_\mathcal{H}$ and
$\mathcal{T}_H$. When the coarse mesh is fixed and $\mathcal{M}\leq 9$,
the best choice of $\mathcal{T}_H$ lies between $\mathcal{T}_\mathcal{H}$
and $\mathcal{T}_h$, demonstrating the advantage of the three-grid ASM over
the two-grid alternative. However, the best overall three-grid configuration
($\mathcal{T}_\mathcal{H}=\mathcal{C}_9$, $\mathcal{T}_H=\mathcal{C}_{10}$)
is only $1.08\times$ faster than the best overall two-grid variant
($\mathcal{T}_\mathcal{H}=\mathcal{T}_H=\mathcal{S}_{9}$).
This is because the coarse \verb|CUDSS| solver performs very well even
for large $\mathcal{T}_\mathcal{H}$ (with reasonable setup times -- see
Figure~\ref{fig:coarse-solver-setup-times}), and therefore it is optimal to choose
$\mathcal{T}_\mathcal{H}$ close to $\mathcal{T}_h$, leaving little room for
improvement by introducing an intermediate $\mathcal{T}_H$. We nevertheless
expect the three-grid approach to offer greater potential on larger problems,
where $\mathcal{T}_\mathcal{H}$ cannot be chosen near $\mathcal{T}_h$.
Such problems, however, turned out to be difficult to handle on a single GPU due to memory
constraints and were therefore left outside the scope of this thesis.

% \begin{figure}
%   \centering
%   \includesvg[width=0.4\textwidth]{coarse_solver_setup_times}
%   \caption{Setup time of the coarse solver (including the construction
%   of $\bm{A_C}$ from~$\bm{A}$) as a~function of $\mathcal{T}_\mathcal{H}$.
%   Parameters: $d=2$, $p=1$ and $\mathcal{T}_h=\mathcal{S}_{11}$.}
%   \label{fig:coarse-solver-setup-times}
% \end{figure}

\subsection{Comparison with AmgX}
\label{sec:experiments-amgx-comparison}

We conclude the experimental evaluation by comparing the performance
of our ASM implementation with the algebraic multigrid preconditioners
provided by NVIDIA AmgX (see Section~\ref{sec:amgx-preconditioner}).
To ensure a fair comparison, we use our PCG implementation
with each preconditioner and measure the total solution time required
to reduce the initial Euclidean norm of the residual by a~factor of $10^9$.
The ASM preconditioner is applied in single precision and tested only with
\verb|Inv(float16)| and \verb|Inv()| as local solvers, in accordance with
the results of the previous sections. Since AmgX also benefits from mixed
precision, its preconditioners are applied in single precision as well.

Results for several model problems with different $d$ and $p$ are reported
in Table~\ref{tab:experiment-speedups}. Note that we use coarser $\mathcal{T}_h$
for larger $p$ to make the problem fit in GPU memory. The problems are large
enough that the direct \verb|CUDSS| solver fails during setup with out-of-memory
errors in most cases; the sole exception is $d=2, p=1, \mathcal{T}_h=\mathcal{S}_{11}$,
for which it performs extraordinarily well, with a solution time of $0.107$\,s
after a $17$\,s setup. Note that when using a direct solver is
applicable, it is expected to offer significantly better performance than
an iterative method, as observed here.

For each problem we list the best AmgX variant
(from the candidates selected in Section~\ref{sec:amgx-preconditioner}),
the best ASM configuration overall, and the best two-grid ASM variants (with
$\mathcal{T}_\mathcal{H}=\mathcal{T}_H$ or $\mathcal{T}_\mathcal{H}=\mathcal{T}_h$).
In all ASM cases the winning configurations used \verb|Inv(float16)|
as the local solver. Note that AmgX solution times reported here are smaller than
those in Section~\ref{sec:amgx-preconditioner}. We attribute this improvement to
replacing the AmgX PCG implementation with our optimized version, which additionally
permits applying the AMG preconditioner in single precision.

\begin{table}
  \centering

  \caption{Best solution times obtained for selected model problems and speedups
  relative to the reference AmgX preconditioner. All variants employ
  \texttt{Inv(float16)} as the local solver.}
  \label{tab:experiment-speedups}

  \begin{tabular}{rrcS[table-format=2.1\,M]|lS[table-format=1.3\,s, round-mode=places, round-precision=3]|ccS[table-format=1.3\,s, round-mode=places, round-precision=3]S[table-format=1.3, round-mode=places, round-precision=3]}
\toprule

 &  &  &  & \multicolumn{2}{c|}{{best AmgX configuration found}} & \multicolumn{4}{c}{{best hybrid variant considered}} \\
$d$ & $p$ & $\mathcal{T}_h$ & {DoFs} & {name} & {time} & {$\mathcal{T}_\mathcal{H}$} & {$\mathcal{T}_H$} & {time} & {speedup} \\

\midrule

\multirow[t]{5}{*}{2} & 1 & $\mathcal{S}_{12}$ & 100.7\,{M} & \verb|L1_TRUNC, fp32| & 2.252932\, \si{\second} & $\mathcal{C}_{11}$ & $\mathcal{C}_{12}$ & 1.824603\, \si{\second} & 1.234752 \,\(\times\) \\
 & 2 & $\mathcal{S}_{11}$ & 50.3\,{M} & \verb|L1_TRUNC, fp32| & 1.840814\, \si{\second} & $\mathcal{C}_{10}$ & $\mathcal{C}_{11}$ & 1.557876\, \si{\second} & 1.181618 \,\(\times\) \\
 & 3 & $\mathcal{S}_{10}$ & 21.0\,{M} & \verb|L1_TRUNC, fp32| & 1.507047\, \si{\second} & $\mathcal{C}_{10}$ & $\mathcal{C}_{10}$ & 1.032556\, \si{\second} & 1.459530 \,\(\times\) \\
 & 4 & $\mathcal{S}_{10}$ & 31.5\,{M} & {--} & {--} & $\mathcal{C}_{10}$ & $\mathcal{C}_{10}$ & 2.027756\, \si{\second} & {$\infty\,\times$} \\
 & 5 & $\mathcal{S}_{9}$ & 11.0\,{M} & \verb|L1_TRUNC, fp32| & 1.425345\, \si{\second} & $\mathcal{C}_{9}$ & $\mathcal{C}_{9}$ & 1.084873\, \si{\second} & 1.313836 \,\(\times\) \\
\multirow[t]{5}{*}{3} & 1 & $\mathcal{S}_{7}$ & 50.3\,{M} & \verb|AGGRESSIVE_L1, fp32| & 1.952959\, \si{\second} & $\mathcal{C}_{6}$ & $\mathcal{C}_{7}$ & 1.407376\, \si{\second} & 1.387660 \,\(\times\) \\
 & 2 & $\mathcal{S}_{6}$ & 15.7\,{M} & \verb|AGGRESSIVE_L1, fp32| & 1.943804\, \si{\second} & $\mathcal{C}_{5}$ & $\mathcal{C}_{6}$ & 1.101878\, \si{\second} & 1.764083 \,\(\times\) \\
 & 3 & $\mathcal{S}_{5}$ & 3.9\,{M} & \verb|L1_TRUNC, fp32| & 1.212115\, \si{\second} & $\mathcal{C}_{5}$ & $\mathcal{C}_{5}$ & 0.542874\, \si{\second} & 2.232773 \,\(\times\) \\
 & 4 & $\mathcal{S}_{5}$ & 6.9\,{M} & \verb|L1_TRUNC, fp64| & 8.055962\, \si{\second} & $\mathcal{C}_{5}$ & $\mathcal{C}_{5}$ & 1.638029\, \si{\second} & 4.918082 \,\(\times\) \\
 & 5 & $\mathcal{S}_{4}$ & 1.4\,{M} & \verb|L1_TRUNC, fp32| & 2.137783\, \si{\second} & $\mathcal{C}_{4}$ & $\mathcal{C}_{4}$ & 0.653890\, \si{\second} & 3.269330 \,\(\times\) \\
\bottomrule
\end{tabular}


  % \bigskip

  % \input{tables/experiment_speedups_2.tex}
\end{table}

For 2D problems, the elapsed time of the solve phase of our ASM implementation
is slightly worse that that of AmgX.
In contrast, for 3D problems the ASM preconditioner outperforms AmgX by a
substantial margin, with speedups between $1.261$ and $1.771\times$.
We emphasize that no further tuning of the example AmgX configurations
was performed here; better performance may be possible by adjusting
those parameters. Moreover, AmgX is a~general-purpose library, whereas our ASM
implementation is specifically tailored to the model problem \eqref{div_eq}.
On the other hand, AmgX is a highly optimized library developed by NVIDIA,
while our implementation obviously has some room for further improvement.
Therefore, we believe that the results are
encouraging and demonstrate the potential of the ASM preconditioner
from~\cite{Przykry}.

From the experiments we observe that
in all cases, optimal or near-optimal performance can be achieved by using
two grids rather than three. This observation is consistent
with Section~\ref{sec:experiments-parallelism}, where we found that for the
problem sizes considered it is optimal to choose $\mathcal{T}_\mathcal{H}$
close to $\mathcal{T}_h$, leaving little room for improvement from an
intermediate mesh. However, it is interesting to note that the classical choice
$\mathcal{T}_H=\mathcal{T}_{\mathcal{H}}$ performed poorly in two test cases
($d=3$, $p\in \{1, 2\}$), while the massively parallel approach
$\mathcal{T}_H=\mathcal{T}_h$ delivered timings close to the best ASM configurations
considered across all tested problems.

\section{Conclusions}

The subject of this thesis is the Additive Schwarz Method (ASM) preconditioner
from \cite{Przykry} for the model elliptic problem \eqref{div_eq}. We adapted
the original method to employ exact local solvers and to support higher-degree
finite elements, and we presented in detail an algorithmic realization
tailored for highly parallel architectures. The objective of the thesis
was to provide an efficient single\nobreakdash-GPU implementation of this method and
to evaluate its performance across a range of parameters, comparing the
best results with the algebraic multigrid solvers from the NVIDIA AmgX library.

Our implementation uses the PyTorch framework with
several targeted optimizations
that may be of independent interest. These include a~custom batched dense
matrix-vector product kernel (which is $10-60\times$ faster than the PyTorch
baseline for small matrix sizes) and a BSR matrix-vector product algorithm
that outperforms cuSPARSE by up to $4.76\times$ for small block sizes.

The experiments confirmed the expected convergence behaviour of the preconditioned
ASM operator in the setting of higher-degree elements. They also showed that
appropriately chosen mixed-precision
configuration have a negligible impact on convergence rate
while providing speedups in the total solution time of
up to $1.81\times$ relative to double precision.

A parameter study revealed that for the fixed coarse mesh the best performance
is achieved for small or moderate local problems, which can be solved efficiently
using the explicit inverse dense batch solver. The three-grid
variant can offer advantages over the two\nobreakdash-grid approach, particularly
when the coarse mesh is substantially coarser than the fine mesh.
However, for problems that fit on a~single GPU it is optimal to choose the coarse
mesh as close as possible to the fine one, leaving
little room for improvement by an intermediate mesh. We therefore expect the
three\nobreakdash-grid method to be more beneficial for larger problems
that do not fit on a~single GPU.

Comparison with NVIDIA AmgX demonstrates that our ASM implementation is competitive:
timings are comparable for 2D problems, while our method yields speedups up to $1.78\times$
over AmgX for 3D problems. These findings attest to the potential of the considered ASM
preconditioner and motivate further research in its highly parallel realizations.

We identify several directions for future experimental work. First, the
implementation can be further optimized by replacing PyTorch-level components
with custom CUDA kernels. Second, given the effectiveness of mixed precision,
investigating inexact \emph{coarse} solvers seems a promising avenue.
Third, extending the approach to multi-GPU setting \cite{Yamazaki} would enable larger problem
sizes and is expected to better exploit the potential of the three-grid method.

\printbibliography

\end{document}


%%% Local Variables:
%%% mode: latex
%%% TeX-master: t
%%% coding: latin-2
%%% End:
